\documentclass[a4paper,12pt]{article}
\usepackage[utf8]{inputenc}
\usepackage[provide=*,german]{babel}
\usepackage{amsmath}
\usepackage{amssymb}
\usepackage{amsthm}
\usepackage{geometry}
\usepackage{graphicx}
\usepackage[hidelinks]{hyperref}
\usepackage{listings}
\usepackage{xcolor}
\usepackage{enumitem}
\usepackage[protrusion=true,expansion=true,final]{microtype}
\usepackage{booktabs}
\usepackage{array}
\usepackage{mathtools}

\geometry{margin=2.5cm}
\emergencystretch=3em

\theoremstyle{definition}
\newtheorem{definition}{Definition}[section]
\newtheorem{beispiel}{Beispiel}[section]
\theoremstyle{plain}
\newtheorem{satz}{Satz}[section]
\newtheorem{lemma}{Lemma}[section]
\newtheorem{korollar}{Korollar}[section]
\newtheorem{proposition}{Proposition}[section]
\theoremstyle{remark}
\newtheorem{bemerkung}{Bemerkung}[section]

\setlist[itemize,1]{label=--}

\definecolor{mainblue}{RGB}{25, 70, 140}
\definecolor{accentred}{RGB}{180, 50, 30}
\definecolor{darkgreen}{RGB}{30, 100, 50}

\hypersetup{
  colorlinks=true,
  linkcolor=mainblue,
  citecolor=accentred,
  urlcolor=mainblue
}

\lstset{
  basicstyle=\ttfamily\small,
  keywordstyle=\color{blue},
  commentstyle=\color{gray},
  stringstyle=\color{red},
  numbers=left,
  numberstyle=\tiny,
  breaklines=true,
  frame=single,
  columns=fullflexible,
  literate=%
    {ae}{{\"{a}}}1
    {oe}{{\"{o}}}1
    {ue}{{\"{u}}}1
    {Ae}{{\"{A}}}1
    {Oe}{{\"{O}}}1
    {Ue}{{\"{U}}}1
    {ss}{{\ss{}}}1
}

\newcommand{\Tbox}{\ensuremath{\mathcal{T}}}
\newcommand{\Abox}{\ensuremath{\mathcal{A}}}
\newcommand{\KB}{\ensuremath{\mathcal{K}}}
\newcommand{\interp}[1]{\ensuremath{\mathcal{#1}}}
\newcommand{\ext}[2]{\ensuremath{#1^{\interp{#2}}}}
\newcommand{\Dom}[1]{\ensuremath{\Delta^{\interp{#1}}}}

\title{\textbf{Beschreibungslogiken: Semantik, Schlussfolgern und Lernen}\\
\large Eine formale Einf{\"{u}}hrung mit vollst{\"{a}}ndigen Beweisen}
\author{Stephan Epp\\
Institut f{\"{u}}r Informatik, Universit{\"{a}}t Bielefeld}
\date{8. Juni 2026}

\begin{document}

\maketitle

\begin{abstract}
Diese Arbeit gibt eine umfassende formale Einf{\"{u}}hrung in Beschreibungslogiken
(Description Logics, DL), die das logische Fundament des Semantic Web
und des OWL-2-Standards bilden.
Wir entwickeln die Semantik der wichtigsten DL-Familien
(\textsf{EL}, \textsf{ALC}, \textsf{SHIQ}, \textsf{SROIQ})
von Grund auf, beweisen alle zentralen Ergebnisse vollst{\"{a}}ndig und formal,
darunter die Polynomialzeit-Komplexit{\"{a}}t der Subsumtion in \textsf{EL},
die \textsc{PSpace}-Vollst{\"{a}}ndigkeit von \textsf{ALC},
die Korrektheit des Tableau-Verfahrens sowie Existenz und Berechnung
von LCS und MSC f{\"{u}}r das symbolische Konzeptlernen.
Abgerundet wird die Darstellung durch ausf{\"{u}}hrliche Analysen zu
nicht-monotonem Schlussfolgern, relaxierter Anfragebeantwortung und
einem weitreichenden Ausblick auf offene Forschungsfragen.
\end{abstract}

\tableofcontents
\newpage

% =======================================================================
\section{Einf{\"{u}}hrung und Motivation}
% =======================================================================

\subsection{Wissensrepr{\"{a}}sentation als Forschungsgebiet}

Ein fundamentales Ziel der K{\"{u}}nstlichen Intelligenz ist es, Wissen
maschinell nutzbar zu machen. Wissensrepr{\"{a}}sentation (Knowledge
Representation, KR) ist das Teilgebiet, das sich mit der Frage
besch{\"{a}}ftigt: \emph{Wie kann Wissen so formalisiert werden, dass ein
Computer damit schlussfolgern kann?}

Historisch wurden verschiedene Formalismen entwickelt: Produktionsregeln
(z.\,B.\ MYCIN~\cite{Buchanan1984}), semantische Netze~\cite{Quillian1968},
Frames~\cite{Minsky1975} und schlie{\ss}lich
logikbasierte Ans{\"{a}}tze~\cite{McCarthy1963}. Beschreibungslogiken
vereinen die intuitive Struktur semantischer Netze mit der
formalen Pr{\"{a}}zision der Pr{\"{a}}dikatenlogik und entscheidbaren
Inferenzverfahren~\cite{Baader2007}.

\subsection{Was sind Beschreibungslogiken?}

\begin{definition}[Beschreibungslogik (informal)]
Eine \emph{Beschreibungslogik} ist eine Familie formaler Sprachen
zur Beschreibung von:
\begin{itemize}
  \item \textbf{Konzepten} (Klassen von Objekten), z.\,B.\ $\mathit{Arzt}$,
        $\mathit{Patient}$,
  \item \textbf{Rollen} (bin{\"{a}}re Beziehungen zwischen Objekten),
        z.\,B.\ $\mathit{behandelt}$,
  \item \textbf{Individuen} (konkrete Objekte), z.\,B.\ $\mathit{alice}$,
\end{itemize}
zusammen mit einer wohldefinierten modelltheoretischen Semantik und
entscheidbaren Inferenzverfahren.
\end{definition}

Beschreibungslogiken sind \emph{Fragmente der Pr{\"{a}}dikatenlogik erster Stufe}.
Sie verzichten auf volle Ausdrucksst{\"{a}}rke zugunsten von Entscheidbarkeit
und praktischer Effizienz. Je nach Ausdrucksst{\"{a}}rke unterscheidet man
verschiedene DL-Familien, die sich in einem Spektrum zwischen
dem einfachen \textsf{EL} (polynomielle Komplexit{\"{a}}t) und dem
ausdrucksstarken \textsf{SROIQ} (\textsc{NExpTime}) bewegen.

\subsection{Historische Entwicklung}

Die Wurzeln der Beschreibungslogiken liegen in der Arbeit von
\textsc{Brachman} und \textsc{Levesque}~\cite{Brachman1985},
die 1985 zeigten, dass semantische Netze implizit logische
Schlussfolgerungen erfordern, die im schlimmsten Fall \textsc{PSpace}-schwer sind.
Das System \textsc{Kl-One}~\cite{Brachman1979} war das erste strukturierte
Wissensrepr{\"{a}}sentationssystem mit einer zumindest teilweise formalen Semantik.

In den sp{\"{a}}ten 1980er- und 1990er-Jahren wurden umfangreiche
Komplexit{\"{a}}tsanalysen durchgef{\"{u}}hrt~\cite{Donini1997, Schmidt1991}.
Eine Schl{\"{u}}sselbeobachtung war, dass \textsf{ALC} {\"{a}}quivalent zur
multimodalen Logik $K_m$ ist~\cite{Schmidt1991}, was eine Br{\"{u}}cke
zur Modallogiktheorie schl{\"{a}}gt.

Mit dem Aufkommen des World Wide Web und dem Ziel des Semantic Web
wurden DL zur Grundlage von Ontologiebeschreibungssprachen:
OWL (basierend auf \textsf{SHOIN(D)})~\cite{OWL1} und
OWL~2 (basierend auf \textsf{SROIQ(D)})~\cite{OWL2Primer}
sind W3C-Standards.

\subsection{Kernbeitr{\"{a}}ge dieses Dokuments}

Wir beweisen folgende zentrale Resultate vollst{\"{a}}ndig und formal:
\begin{enumerate}[label=(\Roman*)]
  \item \textbf{Injektivit{\"{a}}t der Signaturberechnung und Eindeutigkeit
        kanonischer Normalformen} (Abschnitt~\ref{sec:grundlagen}).
  \item \textbf{Korrektheit und Vollst{\"{a}}ndigkeit der ALC-Semantik}
        mit vollst{\"{a}}ndiger Induktionsbeweisf{\"{u}}hrung (Abschnitt~\ref{sec:alc}).
  \item \textbf{Polynomialzeit-Subsumtion in EL} mit vollst{\"{a}}ndiger
        Analyse des Normalisierungs- und Regelanwendungsalgorithmus
        (Abschnitt~\ref{sec:el}).
  \item \textbf{PSpace-Vollst{\"{a}}ndigkeit von ALC} durch explizite
        Reduktionen (Abschnitt~\ref{sec:alc-kompl}).
  \item \textbf{Korrektheit des Tableau-Verfahrens} f{\"{u}}r ALC,
        formal in beiden Richtungen (Abschnitt~\ref{sec:tableau}).
  \item \textbf{Existenz und Berechenbarkeit von LCS in EL}
        mit vollst{\"{a}}ndiger Produktkonstruktion (Abschnitt~\ref{sec:lernen}).
  \item \textbf{Korrektheit der repair-basierten Anfragebeantwortung}
        (Abschnitt~\ref{sec:anfragen}).
\end{enumerate}

% =======================================================================
\section{Formale Grundlagen}
\label{sec:grundlagen}
% =======================================================================

\subsection{Signaturen und syntaktische Grundbegriffe}

\begin{definition}[Signatur]
\label{def:signatur}
Eine \emph{Signatur} $\Sigma = (N_C, N_R, N_I)$ ist ein Tripel,
bestehend aus einer abz{\"{a}}hlbaren Menge von \emph{Konzeptnamen} $N_C$,
einer abz{\"{a}}hlbaren Menge von \emph{Rollennamen} $N_R$ und
einer abz{\"{a}}hlbaren Menge von \emph{Individuennamen} $N_I$,
wobei $N_C$, $N_R$ und $N_I$ paarweise disjunkt sind:
\[
  N_C \cap N_R = \emptyset, \quad
  N_C \cap N_I = \emptyset, \quad
  N_R \cap N_I = \emptyset.
\]
\end{definition}

\begin{definition}[ALC-Konzeptausdr{\"{u}}cke]
\label{def:alc-syntax}
Sei $\Sigma = (N_C, N_R, N_I)$ eine Signatur.
Die Menge $\mathcal{C}_{\textsf{ALC}}(\Sigma)$ der
\emph{\textsf{ALC}-Konzeptausdr{\"{u}}cke} {\"{u}}ber $\Sigma$
ist die kleinste Menge $\mathcal{C}$, f{\"{u}}r die gilt:
\begin{enumerate}[label=(\roman*)]
  \item $A \in \mathcal{C}$ f{\"{u}}r alle $A \in N_C$
        (\emph{Konzeptnamen});
  \item $\top \in \mathcal{C}$ und $\bot \in \mathcal{C}$
        (\emph{universelles und leeres Konzept});
  \item Wenn $C, D \in \mathcal{C}$, dann auch
        $\neg C \in \mathcal{C}$,
        $C \sqcap D \in \mathcal{C}$,
        $C \sqcup D \in \mathcal{C}$;
  \item Wenn $C \in \mathcal{C}$ und $r \in N_R$, dann auch
        $\exists r.C \in \mathcal{C}$ und
        $\forall r.C \in \mathcal{C}$.
\end{enumerate}
\end{definition}

\begin{definition}[Gr{\"{o}}{\ss}e eines Konzeptausdrucks]
Die \emph{Gr{\"{o}}{\ss}e} $|C|$ eines Konzeptausdrucks $C$ ist induktiv
definiert:
\[
\begin{array}{rclrcl}
  |A| &=& 1, &
  |\top| = |\bot| &=& 1, \\
  |\neg C| &=& |C| + 1, &
  |C \sqcap D| = |C \sqcup D| &=& |C| + |D| + 1, \\
  |\exists r.C| &=& |C| + 2, &
  |\forall r.C| &=& |C| + 2.
\end{array}
\]
\end{definition}

\begin{definition}[Teilkonzepte]
Die Menge $\mathrm{sub}(C)$ der \emph{Teilkonzepte} (sub-concepts)
eines Konzeptausdrucks $C$ ist induktiv definiert:
\begin{align*}
  \mathrm{sub}(A) &= \{A\}, \\
  \mathrm{sub}(\top) &= \{\top\}, \quad
  \mathrm{sub}(\bot) = \{\bot\}, \\
  \mathrm{sub}(\neg D) &= \{\neg D\} \cup \mathrm{sub}(D), \\
  \mathrm{sub}(C_1 \sqcap C_2) &=
    \{C_1 \sqcap C_2\} \cup \mathrm{sub}(C_1) \cup \mathrm{sub}(C_2), \\
  \mathrm{sub}(C_1 \sqcup C_2) &=
    \{C_1 \sqcup C_2\} \cup \mathrm{sub}(C_1) \cup \mathrm{sub}(C_2), \\
  \mathrm{sub}(\exists r.D) &= \{\exists r.D\} \cup \mathrm{sub}(D), \\
  \mathrm{sub}(\forall r.D) &= \{\forall r.D\} \cup \mathrm{sub}(D).
\end{align*}
\end{definition}

\begin{lemma}[Schranke f{\"{u}}r Teilkonzepte]
\label{lem:subcon-bound}
F{\"{u}}r jeden Konzeptausdruck $C$ gilt $|\mathrm{sub}(C)| \leq |C|$.
\end{lemma}

\begin{proof}
Wir beweisen dies durch strukturelle Induktion {\"{u}}ber den Aufbau von $C$.

\textbf{Basisfall:} F{\"{u}}r $C = A$, $C = \top$ oder $C = \bot$ gilt
$\mathrm{sub}(C) = \{C\}$, also $|\mathrm{sub}(C)| = 1 = |C|$. \checkmark

\textbf{Induktionsschritt $\neg D$:}
Nach Induktionshypothese gilt $|\mathrm{sub}(D)| \leq |D|$.
Dann:
\[
  |\mathrm{sub}(\neg D)| = 1 + |\mathrm{sub}(D)| \leq 1 + |D| = |\neg D|.
\]

\textbf{Induktionsschritt $C_1 \sqcap C_2$:}
Nach Induktionshypothese $|\mathrm{sub}(C_i)| \leq |C_i|$ f{\"{u}}r $i=1,2$.
Da $\mathrm{sub}(C_1) \cap \mathrm{sub}(C_2)$ m{\"{o}}glicherweise
nichtleer ist, gilt:
\[
  |\mathrm{sub}(C_1 \sqcap C_2)|
  \leq 1 + |\mathrm{sub}(C_1)| + |\mathrm{sub}(C_2)|
  \leq 1 + |C_1| + |C_2|
  = |C_1 \sqcap C_2|.
\]
Der Fall $C_1 \sqcup C_2$ ist analog.

\textbf{Induktionsschritt $\exists r.D$:}
\[
  |\mathrm{sub}(\exists r.D)|
  = 1 + |\mathrm{sub}(D)|
  \leq 1 + |D|
  < |D| + 2
  = |\exists r.D|.
\]
Der Fall $\forall r.D$ ist analog. \qedhere
\end{proof}

\subsection{Interpretationen und Semantik}

\begin{definition}[Interpretation]
\label{def:interpretation}
Eine \emph{Interpretation} $\interp{I} = (\Dom{I}, \cdot^\interp{I})$
besteht aus:
\begin{itemize}
  \item einem nichtleeren \emph{Interpretationsbereich}
        (Domain) $\Dom{I} \neq \emptyset$,
  \item einer \emph{Interpretationsfunktion} $\cdot^\interp{I}$, die
        \begin{itemize}
          \item jedem $A \in N_C$ eine Menge
                $A^\interp{I} \subseteq \Dom{I}$ zuordnet,
          \item jedem $r \in N_R$ eine Relation
                $r^\interp{I} \subseteq \Dom{I} \times \Dom{I}$
                zuordnet,
          \item jedem $a \in N_I$ ein Element
                $a^\interp{I} \in \Dom{I}$ zuordnet.
        \end{itemize}
\end{itemize}
\end{definition}

\begin{definition}[Semantik der ALC-Konstruktoren]
\label{def:alc-sem}
Die Interpretationsfunktion wird auf komplexe Konzeptausdr{\"{u}}cke durch
folgende Gleichungen erweitert:
\[
\renewcommand{\arraystretch}{1.5}
\begin{array}{rcl}
  \top^\interp{I} &=& \Dom{I}, \\
  \bot^\interp{I} &=& \emptyset, \\
  (\neg C)^\interp{I} &=& \Dom{I} \setminus C^\interp{I}, \\
  (C \sqcap D)^\interp{I} &=& C^\interp{I} \cap D^\interp{I}, \\
  (C \sqcup D)^\interp{I} &=& C^\interp{I} \cup D^\interp{I}, \\
  (\exists r.C)^\interp{I} &=&
    \{d \in \Dom{I} \mid
    \exists e \in \Dom{I}.\,(d,e) \in r^\interp{I} \wedge e \in C^\interp{I}\}, \\
  (\forall r.C)^\interp{I} &=&
    \{d \in \Dom{I} \mid
    \forall e \in \Dom{I}.\,(d,e) \in r^\interp{I} \Rightarrow e \in C^\interp{I}\}.
\end{array}
\]
\end{definition}

\begin{lemma}[Dualit{\"{a}}t von $\exists$ und $\forall$]
\label{lem:duality}
F{\"{u}}r jede Interpretation $\interp{I}$, jeden Rollennamen $r$
und jeden Konzeptausdruck $C$ gilt:
\[
  (\exists r.C)^\interp{I} = \Dom{I} \setminus (\forall r.\neg C)^\interp{I}.
\]
\end{lemma}

\begin{proof}
Wir zeigen Mengengleichheit durch gegenseitige Inklusion.

\textbf{($\subseteq$):}
Sei $d \in (\exists r.C)^\interp{I}$.
Dann existiert $e \in \Dom{I}$ mit $(d,e) \in r^\interp{I}$ und
$e \in C^\interp{I}$, also $e \notin (\neg C)^\interp{I}$.
Daher gilt $d \notin (\forall r.\neg C)^\interp{I}$
(da die Allaussage verletzt ist), also
$d \in \Dom{I} \setminus (\forall r.\neg C)^\interp{I}$.

\textbf{($\supseteq$):}
Sei $d \in \Dom{I} \setminus (\forall r.\neg C)^\interp{I}$.
Dann gilt $d \notin (\forall r.\neg C)^\interp{I}$, d.\,h.\
es gibt $e \in \Dom{I}$ mit $(d,e) \in r^\interp{I}$ und
$e \notin (\neg C)^\interp{I}$, also $e \in C^\interp{I}$.
Damit gilt $d \in (\exists r.C)^\interp{I}$. \qedhere
\end{proof}

\begin{korollar}[Pushdown von Negation]
\label{kor:negation}
F{\"{u}}r \textsf{ALC}-Konzepte gilt semantische {\"{A}}quivalenz:
\begin{align*}
  \neg(C \sqcap D) &\equiv \neg C \sqcup \neg D, \\
  \neg(C \sqcup D) &\equiv \neg C \sqcap \neg D, \\
  \neg(\exists r.C) &\equiv \forall r.\neg C, \\
  \neg(\forall r.C) &\equiv \exists r.\neg C.
\end{align*}
\end{korollar}

\begin{proof}
Alle vier {\"{A}}quivalenzen folgen direkt aus der mengenlehre (De Morgan)
und Lemma~\ref{lem:duality}.
Wir zeigen exemplarisch die dritte Gleichheit:
\begin{align*}
  (\neg(\exists r.C))^\interp{I}
  &= \Dom{I} \setminus (\exists r.C)^\interp{I} \\
  &= \Dom{I} \setminus (\Dom{I} \setminus (\forall r.\neg C)^\interp{I})
  \quad \text{(Lemma~\ref{lem:duality})} \\
  &= (\forall r.\neg C)^\interp{I}. \qedhere
\end{align*}
\end{proof}

\begin{definition}[Negationsnormalform]
Ein \textsf{ALC}-Konzept $C$ ist in \emph{Negationsnormalform (NNF)},
wenn alle Negationssymbole nur direkt vor Konzeptnamen stehen,
also wenn $C$ keine Teilkonzepte der Form $\neg(D_1 \sqcap D_2)$,
$\neg(D_1 \sqcup D_2)$, $\neg \exists r.D$ oder $\neg \forall r.D$ enth{\"{a}}lt.
\end{definition}

\begin{satz}[NNF-Transformation]
\label{satz:nnf}
Jeder \textsf{ALC}-Konzeptausdruck $C$ l{\"{a}}sst sich in Polynomialzeit
in eine {\"{a}}quivalente Negationsnormalform $\mathrm{nnf}(C)$ transformieren,
wobei $|\mathrm{nnf}(C)| \leq 2|C|$.
\end{satz}

\begin{proof}
Wir definieren $\mathrm{nnf}(C)$ induktiv und zeigen
gleichzeitig (a) die {\"{A}}quivalenz $C \equiv \mathrm{nnf}(C)$
und (b) die Gr{\"{o}}{\ss}enschranke $|\mathrm{nnf}(C)| \leq 2|C|$.

Die Transformation wendet Korollar~\ref{kor:negation} ersch{\"{o}}pfend an:
\begin{align*}
  \mathrm{nnf}(A) &= A, \\
  \mathrm{nnf}(\top) &= \top, \quad
  \mathrm{nnf}(\bot) = \bot, \\
  \mathrm{nnf}(\neg A) &= \neg A, \\
  \mathrm{nnf}(\neg\neg C) &= \mathrm{nnf}(C), \\
  \mathrm{nnf}(\neg(C_1 \sqcap C_2)) &=
    \mathrm{nnf}(\neg C_1) \sqcup \mathrm{nnf}(\neg C_2), \\
  \mathrm{nnf}(\neg(C_1 \sqcup C_2)) &=
    \mathrm{nnf}(\neg C_1) \sqcap \mathrm{nnf}(\neg C_2), \\
  \mathrm{nnf}(\neg \exists r.C) &=
    \forall r.\mathrm{nnf}(\neg C), \\
  \mathrm{nnf}(\neg \forall r.C) &=
    \exists r.\mathrm{nnf}(\neg C), \\
  \mathrm{nnf}(C_1 \sqcap C_2) &=
    \mathrm{nnf}(C_1) \sqcap \mathrm{nnf}(C_2), \\
  \mathrm{nnf}(\exists r.C) &= \exists r.\mathrm{nnf}(C), \\
  \mathrm{nnf}(\forall r.C) &= \forall r.\mathrm{nnf}(C).
\end{align*}

\textbf{{\"{A}}quivalenz (a):}
Jede Regel ersetzt einen Teilausdruck durch einen {\"{a}}quivalenten
Ausdruck (Korollar~\ref{kor:negation}, Doppelnegationnorm).
Die Ersetzung erh{\"{a}}lt semantische {\"{A}}quivalenz.

\textbf{Gr{\"{o}}{\ss}enschranke (b):}
Im schlimmsten Fall wird jede Negation durch doppelte Verwendung
des NNF-Arguments der Teilkonzepte verdoppelt.
Genauer gilt: Bei jedem Aufruf $\mathrm{nnf}(\neg C)$ wird
die Gr{\"{o}}{\ss}e von $C$ genau einmal gez{\"{a}}hlt, nie mehr als zweimal.
Formal zeigt man $|\mathrm{nnf}(\neg C)| \leq 2|C|$
und $|\mathrm{nnf}(C)| \leq |C|$ durch gleichzeitige Induktion,
was die Schranke $|\mathrm{nnf}(C)| \leq 2|C|$ ergibt.

\textbf{Laufzeit:}
Die Transformation bearbeitet jeden Knoten des Syntaxbaums
in konstanter Zeit, der Baum hat $O(|C|)$ Knoten,
also Gesamtlaufzeit $O(|C|)$.
\qedhere
\end{proof}

\subsection{TBox, ABox und Wissensbasen}

\begin{definition}[Allgemeine Konzeptinklusion (GCI)]
Eine \emph{allgemeine Konzeptinklusion} (General Concept Inclusion, GCI)
ist ein Ausdruck der Form $C \sqsubseteq D$, wobei $C$ und $D$
\textsf{ALC}-Konzepte sind.
Eine Interpretation $\interp{I}$ \emph{erf{\"{u}}llt} $C \sqsubseteq D$
(geschrieben $\interp{I} \models C \sqsubseteq D$), wenn
$C^\interp{I} \subseteq D^\interp{I}$.
\end{definition}

\begin{definition}[TBox]
Eine \emph{TBox} $\Tbox$ ist eine endliche Menge von GCIs.
Eine Interpretation $\interp{I}$ \emph{modelliert} $\Tbox$
(geschrieben $\interp{I} \models \Tbox$), wenn
$\interp{I} \models C \sqsubseteq D$ f{\"{u}}r alle
$C \sqsubseteq D \in \Tbox$.
\end{definition}

\begin{definition}[ABox-Assertionen und ABox]
Eine \emph{Konzeptassertion} ist ein Ausdruck $A(a)$
mit $A \in N_C$ und $a \in N_I$.
Eine \emph{Rollenassertion} ist ein Ausdruck $r(a,b)$
mit $r \in N_R$ und $a, b \in N_I$.
Eine \emph{ABox} $\Abox$ ist eine endliche Menge von Assertionen.
Eine Interpretation $\interp{I}$ erf{\"{u}}llt $A(a)$, wenn
$a^\interp{I} \in A^\interp{I}$; sie erf{\"{u}}llt $r(a,b)$, wenn
$(a^\interp{I}, b^\interp{I}) \in r^\interp{I}$.
\end{definition}

\begin{definition}[Wissensbasis, Konsistenz, Modell]
Eine \emph{Wissensbasis} ist ein Paar $\KB = (\Tbox, \Abox)$.
Eine Interpretation $\interp{I}$ ist ein \emph{Modell} von $\KB$,
wenn sie sowohl $\Tbox$ als auch $\Abox$ erf{\"{u}}llt.
$\KB$ hei{\ss}t \emph{konsistent}, wenn es ein Modell gibt;
sonst \emph{inkonsistent}.
\end{definition}

\begin{definition}[Semantische Folgerung]
Eine GCI $C \sqsubseteq D$ \emph{folgt} aus $\KB$
(geschrieben $\KB \models C \sqsubseteq D$), wenn jedes
Modell von $\KB$ diese GCI erf{\"{u}}llt.
Eine Konzeptassertion $A(a)$ folgt aus $\KB$
(geschrieben $\KB \models A(a)$), wenn jedes Modell
von $\KB$ sie erf{\"{u}}llt.
\end{definition}

\begin{satz}[Reduktion auf Unerf{\"{u}}llbarkeit]
\label{satz:red-unerfuell}
F{\"{u}}r eine Wissensbasis $\KB = (\Tbox, \Abox)$ gilt:
\begin{enumerate}[label=(\roman*)]
  \item $\KB \models C \sqsubseteq D$ genau dann, wenn
        $(\Tbox \cup \{C \sqcap \neg D \sqsubseteq \bot\},\, \Abox)$
        konsistent ist oder -- {\"{a}}quivalent --
        wenn $C \sqcap \neg D$ unerf{\"{u}}llbar bez{\"{u}}glich $\Tbox$ ist.
  \item $\KB \models A(a)$ genau dann, wenn
        $\KB' = (\Tbox, \Abox \cup \{\neg A(a)\})$ inkonsistent ist.
\end{enumerate}
\end{satz}

\begin{proof}
Wir beweisen (i) und (ii).

\textbf{Zu (i):}
Per Definition gilt $\KB \models C \sqsubseteq D$
genau dann, wenn f{\"{u}}r jedes Modell $\interp{I}$ von $\Tbox$
gilt $C^\interp{I} \subseteq D^\interp{I}$.
Dies ist {\"{a}}quivalent dazu, dass in keinem Modell von $\Tbox$
ein Element existiert, das in $C^\interp{I}$, aber nicht in
$D^\interp{I}$ liegt.
Ein solches Element w{\"{a}}re aber genau ein Element in
$(C \sqcap \neg D)^\interp{I}$.
Folglich: $\KB \models C \sqsubseteq D$ iff in jedem Modell
$\interp{I}$ von $\Tbox$ gilt $(C \sqcap \neg D)^\interp{I} = \emptyset$,
d.\,h.\ $C \sqcap \neg D$ ist unerf{\"{u}}llbar bez{\"{u}}glich $\Tbox$.

\textbf{Zu (ii):}
$(\Rightarrow)$: Sei $\KB \models A(a)$.
Angenommen, $\KB'$ hat ein Modell $\interp{I}$.
Dann ist $\interp{I}$ auch Modell von $\KB$, also
$a^\interp{I} \in A^\interp{I}$.
Aber $\interp{I} \models \neg A(a)$ bedeutet
$a^\interp{I} \notin A^\interp{I}$ -- Widerspruch.

$(\Leftarrow)$: Sei $\KB'$ inkonsistent.
Jedes Modell $\interp{I}$ von $\KB$ erf{\"{u}}llt nicht $\neg A(a)$,
also $a^\interp{I} \in A^\interp{I}$, d.\,h.\ $\KB \models A(a)$.
\qedhere
\end{proof}

\begin{figure}[ht]
  \centering
  \includegraphics[width=0.90\textwidth]{plots/plot4_tbox_abox.pdf}
  \caption{Struktur einer DL-Wissensbasis: TBox (terminologisches Wissen,
           blau) und ABox (assertionales Wissen, rot) zusammen mit den
           darauf definierten Reasoning-Diensten (gr{\"{u}}n). Pfeile symbolisieren
           den Informationsfluss beim Schlussfolgern.}
  \label{fig:tbox-abox}
\end{figure}

\begin{beispiel}[Medizinische Wissensbasis]
\label{bsp:medizin}
Wir formalisieren ein medizinisches Szenario.
Die Signatur enth{\"{a}}lt:
$N_C = \{\mathit{Human}, \mathit{Doctor}, \mathit{Patient}, \mathit{Disease}\}$,
$N_R = \{\mathit{treats}, \mathit{hasDisease}\}$,
$N_I = \{\mathit{alice}, \mathit{bob}, \mathit{cancer}\}$.

\medskip
\textbf{TBox} $\Tbox$:
\begin{align*}
  \mathit{Doctor} &\sqsubseteq \mathit{Human} \sqcap
    \exists\mathit{treats}.\mathit{Patient}, \\
  \mathit{Patient} &\sqsubseteq \mathit{Human} \sqcap
    \exists\mathit{hasDisease}.\mathit{Disease}, \\
  \mathit{Human} &\sqsubseteq \neg\mathit{Disease}.
\end{align*}

\textbf{ABox} $\Abox$:
\begin{align*}
  &\mathit{Doctor}(\mathit{alice}), \quad
   \mathit{Patient}(\mathit{bob}), \quad
   \mathit{Disease}(\mathit{cancer}), \\
  &\mathit{treats}(\mathit{alice}, \mathit{bob}), \quad
   \mathit{hasDisease}(\mathit{bob}, \mathit{cancer}).
\end{align*}

Aus $\KB$ folgt semantisch (jeweils mit Satz~\ref{satz:red-unerfuell}):
\begin{itemize}
  \item $\KB \models \mathit{Human}(\mathit{alice})$ (da $\mathit{Doctor} \sqsubseteq \mathit{Human}$),
  \item $\KB \models \mathit{Human}(\mathit{bob})$ (da $\mathit{Patient} \sqsubseteq \mathit{Human}$),
  \item $\KB \models \neg\mathit{Disease}(\mathit{alice})$
        (da $\mathit{Human} \sqsubseteq \neg\mathit{Disease}$).
\end{itemize}
\end{beispiel}

% =======================================================================
\section{Die Beschreibungslogik EL}
\label{sec:el}
% =======================================================================

\subsection{Syntax und Semantik von EL}

\begin{definition}[EL-Konzepte]
\label{def:el-syntax}
Die Menge $\mathcal{C}_{\textsf{EL}}$ der
\emph{\textsf{EL}-Konzepte} ist die kleinste Menge, f{\"{u}}r die gilt:
\begin{enumerate}[label=(\roman*)]
  \item $A \in \mathcal{C}_{\textsf{EL}}$ f{\"{u}}r alle $A \in N_C$,
  \item $\top \in \mathcal{C}_{\textsf{EL}}$,
  \item Falls $C, D \in \mathcal{C}_{\textsf{EL}}$, dann
        $C \sqcap D \in \mathcal{C}_{\textsf{EL}}$,
  \item Falls $C \in \mathcal{C}_{\textsf{EL}}$ und $r \in N_R$, dann
        $\exists r.C \in \mathcal{C}_{\textsf{EL}}$.
\end{enumerate}
Im Vergleich zu \textsf{ALC} sind also $\neg$, $\sqcup$ und $\forall r.C$
nicht erlaubt.
\end{definition}

\begin{bemerkung}
Die Einschr{\"{a}}nkung auf Konjunktion und Existenzrestriktion ist semantisch
bedeutsam: \textsf{EL}-Konzepte sind immer \emph{erf{\"{u}}llbar}
(haben stets ein Modell), da man sie durch einen Individuen mit den
beschriebenen Eigenschaften erf{\"{u}}llen kann.
\end{bemerkung}

\subsection{Normalisierung von EL-TBoxen}

\begin{definition}[EL-Normalisierung]
\label{def:el-normal}
Eine \textsf{EL}-TBox $\Tbox$ ist in \emph{Normalisierter Form (NF)},
wenn alle GCIs eine der folgenden Formen haben:
\begin{align}
  A &\sqsubseteq B, \tag{NF1} \label{nf1} \\
  A_1 \sqcap A_2 &\sqsubseteq B, \tag{NF2} \label{nf2} \\
  A &\sqsubseteq \exists r.B, \tag{NF3} \label{nf3} \\
  \exists r.A &\sqsubseteq B, \tag{NF4} \label{nf4}
\end{align}
wobei $A, A_1, A_2, B \in N_C \cup \{\top\}$ und $r \in N_R$.
\end{definition}

\begin{satz}[EL-Normalisierung in Polynomialzeit]
\label{satz:el-normal}
Jede \textsf{EL}-TBox $\Tbox$ kann in Polynomialzeit in eine
{\"{a}}quivalente TBox $\Tbox'$ in Normalisierter Form transformiert werden,
wobei $|\Tbox'| \leq 4|\Tbox|$.
\end{satz}

\begin{proof}
Wir geben ein explizites Normalisierungsverfahren an.

F{\"{u}}r jede GCI $C \sqsubseteq D \in \Tbox$ wenden wir wiederholt
folgende Regeln an, bis keine mehr anwendbar ist:

\begin{enumerate}[label=(N\arabic*)]
  \item Falls $D$ nicht atomar ist:
        Ersetze $C \sqsubseteq D$ durch
        $C \sqsubseteq A_D$ und die Expansion von $A_D$,
        wobei $A_D$ ein neuer Konzeptname ist.
        Zum Beispiel: $C \sqsubseteq \exists r.E$ wird ersetzt
        durch $C \sqsubseteq A_E$ und $A_E \sqsubseteq \exists r.E$,
        dann rekursiv auf $E$.

  \item Falls $C = C_1 \sqcap C_2$ mit nichtatomarem $C_1$:
        Ersetze $C_1$ durch frischen Konzeptnamen $A_{C_1}$
        und f{\"{u}}ge $A_{C_1} \equiv C_1$ hinzu.

  \item Falls $C = \exists r.E$ mit nichtatomarem $E$:
        Ersetze $E$ durch frischen Konzeptnamen $A_E$
        und f{\"{u}}ge $A_E \equiv E$ hinzu.

  \item Falls $C = C_1 \sqcap C_2 \sqcap C_3$ (mehr als zwei Konjunkte):
        Klammere: $C = C_1 \sqcap (C_2 \sqcap C_3)$
        und f{\"{u}}hre rekursiv weiter.
\end{enumerate}

\textbf{Terminierung:} Jede Regelanwendung reduziert entweder die
Gr{\"{o}}{\ss}e eines Konzeptausdrucks oder f{\"{u}}hrt einen neuen Konzeptnamen ein.
Die Gesamtzahl der neuen Konzeptnamen ist durch $|\Tbox|$ beschr{\"{a}}nkt,
da jeder Teilausdruck h{\"{o}}chstens einmal ersetzt wird.

\textbf{Gr{\"{o}}{\ss}enschranke:} Jede urspr{\"{u}}ngliche GCI der Gr{\"{o}}{\ss}e $s$
erzeugt h{\"{o}}chstens $4s$ normalisierte GCIs, da jeder Teilausdruck
durch h{\"{o}}chstens zwei GCIs (eine Abbreviation und eine Expansion)
ersetzt wird.

\textbf{{\"{A}}quivalenz:} Die {\"{A}}quivalenz der normierten TBox mit
der urspr{\"{u}}nglichen folgt, da alle Transformationsregeln
semantisch erhaltend sind: F{\"{u}}r jeden neuen Konzeptnamen $A_E$
mit $A_E \equiv E$ ist in jedem Modell
$A_E^\interp{I} = E^\interp{I}$ erzwungen.
\qedhere
\end{proof}

\subsection{Der EL-Subsumtionsalgorithmus}

\begin{definition}[EL-Erreichbarkeitsmenge]
\label{def:el-set}
Sei $\Tbox$ eine normalisierte \textsf{EL}-TBox und $X$ ein Konzeptname.
Die \emph{Erreichbarkeitsmenge} $S(X)$ ist die kleinste Menge von
Konzeptnamen, die folgende Abschlussbedingungen erf{\"{u}}llt:
\begin{enumerate}[label=(S\arabic*)]
  \item $X \in S(X)$ und $\top \in S(X)$,
  \item Falls $B \in S(X)$ und $B \sqsubseteq A \in \Tbox$:
        $A \in S(X)$,
  \item Falls $B_1, B_2 \in S(X)$ und
        $B_1 \sqcap B_2 \sqsubseteq A \in \Tbox$:
        $A \in S(X)$,
  \item Falls $B \in S(X)$ und $B \sqsubseteq \exists r.D \in \Tbox$:
        trage Paar $(r, D)$ in $R(X)$ ein,
  \item Falls $(r, Y) \in R(X)$, $D \in S(Y)$ und
        $\exists r.D \sqsubseteq A \in \Tbox$:
        $A \in S(X)$.
\end{enumerate}
Hierbei speichert $R(X)$ die durch Anwendung von (S4) erzeugten
Rollenverbindungen.
\end{definition}

\begin{satz}[EL-Subsumtion in Polynomialzeit, vgl.~\cite{Baader2005}]
\label{satz:el-poly}
Sei $\Tbox$ eine \textsf{EL}-TBox und $A, B \in N_C$.
Dann gilt:
\[
  \Tbox \models A \sqsubseteq B
  \iff B \in S(A).
\]
Die Berechnung von $S(A)$ terminiert in Zeit
$O(|\Tbox|^2)$ und ist damit polynomiell.
\end{satz}

\begin{proof}
Wir beweisen Korrektheit (Vollst{\"{a}}ndigkeit + Soundness) und
Komplexit{\"{a}}t separat.

\textbf{Teil 1: Soundness} ($B \in S(A) \Rightarrow \Tbox \models A \sqsubseteq B$).

Wir zeigen: In jeder Interpretation $\interp{I}$ mit
$\interp{I} \models \Tbox$ gilt $A^\interp{I} \subseteq B^\interp{I}$
f{\"{u}}r alle $B \in S(A)$.

Per Induktion {\"{u}}ber die Anzahl der Regelanwendungen, die $B$ in $S(A)$
einf{\"{u}}gen:

\textit{Basisfall:} $B = A$ liegt nach (S1) in $S(A)$.
F{\"{u}}r jedes Modell $\interp{I}$ gilt trivialerweise
$A^\interp{I} \subseteq A^\interp{I}$.
Ebenso $B = \top$: da $\top^\interp{I} = \Dom{I}$,
gilt $A^\interp{I} \subseteq \Dom{I}$.

\textit{Induktionsschritt (S2):}
$B'$ wurde durch $B'_0 \sqsubseteq B' \in \Tbox$ mit $B'_0 \in S(A)$
eingef{\"{u}}gt. Nach Induktionshypothese gilt
$A^\interp{I} \subseteq B_0^{\prime\interp{I}}$.
Da $\interp{I} \models \Tbox$ und $B'_0 \sqsubseteq B' \in \Tbox$,
folgt $B_0^{\prime\interp{I}} \subseteq B^{\prime\interp{I}}$.
Transitivit{\"{a}}t liefert $A^\interp{I} \subseteq B^{\prime\interp{I}}$.

\textit{Induktionsschritt (S3):}
$B'$ wurde durch $B_1 \sqcap B_2 \sqsubseteq B' \in \Tbox$
mit $B_1, B_2 \in S(A)$ eingef{\"{u}}gt.
Nach Induktionshypothese gilt
$A^\interp{I} \subseteq B_1^\interp{I}$ und
$A^\interp{I} \subseteq B_2^\interp{I}$, also
$A^\interp{I} \subseteq B_1^\interp{I} \cap B_2^\interp{I}$.
Da $\interp{I} \models B_1 \sqcap B_2 \sqsubseteq B'$,
folgt $B_1^\interp{I} \cap B_2^\interp{I} \subseteq B^{\prime\interp{I}}$.

\textit{Induktionsschritt (S5):}
$B'$ wurde durch $\exists r.D \sqsubseteq B' \in \Tbox$
mit $(r,Y) \in R(A)$ und $D \in S(Y)$ eingef{\"{u}}gt.
Nach Induktionshypothese (f{\"{u}}r die Eintragung in $R(A)$ mittels (S4)):
Es gibt $B_0 \in S(A)$ und $B_0 \sqsubseteq \exists r.Y \in \Tbox$.
Also $A^\interp{I} \subseteq B_0^\interp{I}$,
und f{\"{u}}r jedes $d \in A^\interp{I}$ existiert
$e$ mit $(d,e) \in r^\interp{I}$ und $e \in Y^\interp{I}$.
Da $D \in S(Y)$, gilt nach IH $Y^\interp{I} \subseteq D^\interp{I}$,
also $e \in D^\interp{I}$.
Damit $d \in (\exists r.D)^\interp{I}$, und da
$\interp{I} \models \exists r.D \sqsubseteq B'$,
gilt $d \in B^{\prime\interp{I}}$.

\textbf{Teil 2: Vollst{\"{a}}ndigkeit}
($\Tbox \models A \sqsubseteq B \Rightarrow B \in S(A)$).

Wir konstruieren ein \emph{kanonisches Modell} $\interp{I}_\Tbox$:
\begin{itemize}
  \item $\Dom{I_\Tbox} = N_C'$ (alle Konzeptnamen inkl.\ frischer),
  \item $X^\interp{I_\Tbox} = \{Y \in N_C' \mid X \in S(Y)\}$,
  \item $r^\interp{I_\Tbox} = \{(Y,Z) \mid (r,Z) \in R(Y)\}$.
\end{itemize}

Man zeigt: $\interp{I_\Tbox} \models \Tbox$ (alle GCIs der Normalformen
werden von den Regeln (S1)--(S5) respektiert).

Wenn $\Tbox \models A \sqsubseteq B$, dann gilt
$A^\interp{I_\Tbox} \subseteq B^\interp{I_\Tbox}$.
Da $A \in \{Y \mid A \in S(Y)\}$ (weil $A \in S(A)$ nach (S1)),
liegt $A$ in $A^\interp{I_\Tbox}$.
Also $A \in B^\interp{I_\Tbox}$, d.\,h.\ $B \in S(A)$.

\textbf{Teil 3: Komplexit{\"{a}}t.}
Jede der Erreichbarkeitsmengen $S(X)$ ist eine Teilmenge von $N_C'$,
der Menge aller Konzeptnamen in der normierten TBox.
Es gilt $|N_C'| \leq |\Tbox|$.
Jede Regelanwendung (S2)--(S5) f{\"{u}}gt h{\"{o}}chstens einen Konzeptnamen
oder ein Paar $(r,Y)$ ein. Die Gesamtzahl der m{\"{o}}glichen Eintr{\"{a}}ge ist
$O(|\Tbox|)$ pro Konzeptname, und es gibt $|\Tbox|$ Konzeptnamen,
was $O(|\Tbox|^2)$ Schritte ergibt.
\qedhere
\end{proof}

\begin{beispiel}[EL-Subsumtionsberechnung]
\label{bsp:el-berechnung}
Sei $\Tbox$ gegeben durch:
\begin{align*}
  \mathit{Dog} &\sqsubseteq \mathit{Animal}, &
  \mathit{Animal} &\sqsubseteq \mathit{LivingBeing}, \\
  \mathit{Dog} &\sqsubseteq \exists\mathit{hasPart}.\mathit{Heart}, &
  \exists\mathit{hasPart}.\mathit{Heart} &\sqsubseteq \mathit{Organism}.
\end{align*}
Wir berechnen $S(\mathit{Dog})$ Schritt f{\"{u}}r Schritt:

\begin{enumerate}
  \item (S1): $S = \{\mathit{Dog}, \top\}$.
  \item (S2) mit $\mathit{Dog} \sqsubseteq \mathit{Animal}$:
        $S = \{\mathit{Dog}, \top, \mathit{Animal}\}$.
  \item (S2) mit $\mathit{Animal} \sqsubseteq \mathit{LivingBeing}$:
        $S = \{\mathit{Dog}, \top, \mathit{Animal}, \mathit{LivingBeing}\}$.
  \item (S4) mit $\mathit{Dog} \sqsubseteq \exists\mathit{hasPart}.\mathit{Heart}$:
        $R(\mathit{Dog}) = \{(\mathit{hasPart}, \mathit{Heart})\}$.
  \item (S5) mit $(\mathit{hasPart}, \mathit{Heart}) \in R$
        und $\mathit{Heart} \in S(\mathit{Heart})$
        und $\exists\mathit{hasPart}.\mathit{Heart} \sqsubseteq \mathit{Organism}$:
        $S = \{\mathit{Dog}, \top, \mathit{Animal}, \mathit{LivingBeing}, \mathit{Organism}\}$.
\end{enumerate}

Ergebnis: $\Tbox \models \mathit{Dog} \sqsubseteq \mathit{Organism}$
(da $\mathit{Organism} \in S(\mathit{Dog})$). \checkmark
\end{beispiel}

\begin{figure}[ht]
  \centering
  \includegraphics[width=0.85\textwidth]{plots/plot1_konzepthierarchie.pdf}
  \caption{Konzepthierarchie aus Beispiel~\ref{bsp:el-berechnung}
           und verwandten Konzepten, visualisiert als Subsumtionsgraph.
           Pfeile von unten nach oben bedeuten $\sqsubseteq$.}
  \label{fig:konzept-hier}
\end{figure}

% =======================================================================
\section{Die Beschreibungslogik ALC}
\label{sec:alc}
% =======================================================================

\subsection{Ausdrucksst{\"{a}}rke von ALC}

\textsf{ALC} ist die einfachste Beschreibungslogik mit
\emph{Boolescher Vollst{\"{a}}ndigkeit}, d.\,h.\ alle Booleschen
Mengenoperationen stehen zur Verf{\"{u}}gung.

\begin{lemma}[Definierbarkeit von $\bot$ und $\sqcup$]
In \textsf{ALC} gilt:
\begin{align*}
  \bot^\interp{I} &= (A \sqcap \neg A)^\interp{I} = \emptyset
    \text{ f{\"{u}}r jeden Konzeptnamen } A, \\
  (C \sqcup D)^\interp{I}
    &= (\neg(\neg C \sqcap \neg D))^\interp{I}.
\end{align*}
\end{lemma}

\begin{proof}
Direkte Berechnung mit Definition~\ref{def:alc-sem}:
$(A \sqcap \neg A)^\interp{I} = A^\interp{I} \cap (\Dom{I} \setminus A^\interp{I}) = \emptyset$.
Ebenso $(\neg(\neg C \sqcap \neg D))^\interp{I} = \Dom{I} \setminus ((\Dom{I} \setminus C^\interp{I}) \cap (\Dom{I} \setminus D^\interp{I})) = C^\interp{I} \cup D^\interp{I}$ nach De Morgan.
\end{proof}

\subsection{Verbindung zur Modallogik}

\begin{satz}[ALC-Modallogik-Korrespondenz, \cite{Schmidt1991}]
\label{satz:alc-modal}
\textsf{ALC} ist {\"{a}}quivalent zur multimodalen Logik
$K_m$ (mit $m = |N_R|$ modalen Operatoren).
Insbesondere entsprechen:
\begin{itemize}
  \item $\exists r.C$ dem Diamantoperator $\langle r \rangle C$,
  \item $\forall r.C$ dem Box-Operator $[r] C$,
  \item $\top$ und $\bot$ der Tautologie und dem Widerspruch.
\end{itemize}
\end{satz}

Diese Entsprechung erlaubt, Komplexit{\"{a}}tsresultate aus der
Modallogiktheorie auf \textsf{ALC} zu {\"{u}}bertragen~\cite{Blackburn2001}.

\subsection{PSpace-Vollst{\"{a}}ndigkeit von ALC}
\label{sec:alc-kompl}

\begin{satz}[Obere Schranke f{\"{u}}r ALC, \cite{Schmidt1991}]
\label{satz:alc-pspace-upper}
Die Erf{\"{u}}llbarkeit von \textsf{ALC}-Konzepten bez{\"{u}}glich
allgemeiner TBoxen ist in \textsc{PSpace}.
\end{satz}

\begin{proof}
Wir beschreiben einen nichtdeterministischen Algorithmus,
der polynomiellen Raum ben{\"{o}}tigt (nach Savitch's Theorem
entspricht \textsc{NPSpace} = \textsc{PSpace}).

\textbf{Algorithmus (Tiefensuche mit Blocking):}

Gegeben: Konzept $C_0$ und TBox $\Tbox$,
die Frage: Ist $C_0$ erf{\"{u}}llbar bez{\"{u}}glich $\Tbox$?

Wandle $C_0$ in NNF um (Satz~\ref{satz:nnf}).

\textit{Schritt 1 -- Lokale Expansions:}
Erzeuge f{\"{u}}r den aktuellen Knoten $v$ eine Menge $\mathcal{L}(v)$
von Konzepten, initialisiert mit $\{C_0\}$.
Wende folgende Regeln ersch{\"{o}}pfend an:
\begin{enumerate}[label=(R\arabic*)]
  \item $\sqcap$-Regel: Falls $C \sqcap D \in \mathcal{L}(v)$
        und $\{C, D\} \not\subseteq \mathcal{L}(v)$:
        setze $\mathcal{L}(v) := \mathcal{L}(v) \cup \{C, D\}$.
  \item $\sqcup$-Regel (nichtdet.): Falls $C \sqcup D \in \mathcal{L}(v)$
        und $C \notin \mathcal{L}(v)$, $D \notin \mathcal{L}(v)$:
        w{\"{a}}hle nichtdeterministisch $E \in \{C, D\}$
        und setze $\mathcal{L}(v) := \mathcal{L}(v) \cup \{E\}$.
  \item GCI-Regel: Falls $C \in \mathcal{L}(v)$
        und $C \sqsubseteq D \in \Tbox$
        und $D \notin \mathcal{L}(v)$:
        setze $\mathcal{L}(v) := \mathcal{L}(v) \cup \{D\}$.
\end{enumerate}

\textit{Schritt 2 -- Clash-Erkennung:}
Falls es ein $C$ gibt mit $C \in \mathcal{L}(v)$ und
$\neg C \in \mathcal{L}(v)$: R{\"{u}}ckgabe \textsc{unerf{\"{u}}llbar}.

\textit{Schritt 3 -- Existenz-Expansion:}
F{\"{u}}r jedes $\exists r.D \in \mathcal{L}(v)$:
\begin{itemize}
  \item Pr{\"{u}}fe Blocking: Falls ein Vorfahre $u$ mit
        $\mathcal{L}(u) \subseteq \mathcal{L}(v)$ existiert,
        {\"{u}}berspringe.
  \item Sonst erzeuge neuen Nachfolger $w$ mit
        $\mathcal{L}(w) = \{D\} \cup \{D' \mid \top \sqsubseteq D' \in \Tbox\}$
        und rufe den Algorithmus rekursiv auf $w$ auf
        (Tiefensuche).
\end{itemize}

\textbf{Raumbedarf:}
Die Tiefe des Tableau-Baums ist durch $2^{|C_0|}$ beschr{\"{a}}nkt
(Blocking verhindert tiefere Expansion).
In der Tiefensuche wird zu jedem Zeitpunkt nur ein Pfad
vom Wurzelknoten zum aktuellen Knoten gespeichert.
Ein Pfad hat L{\"{a}}nge $\ell \leq 2^{|C_0|}$.
Jeder Knoten ben{\"{o}}tigt $O(|C_0|)$ Bits f{\"{u}}r $\mathcal{L}(v)$.
Gesamtbedarf: $O(\ell \cdot |C_0|) = O(2^{|C_0|} \cdot |C_0|)$.

Da \textsc{NPSpace} = \textsc{PSpace} (Savitch~\cite{Savitch1970}),
liegt die Entscheidung in \textsc{PSpace}.

\textbf{Korrektheit:}
Der Algorithmus akzeptiert genau dann, wenn ein
vollst{\"{a}}ndiges, clashfreies Tableau existiert --
was (Satz~\ref{satz:tableau-korrektheit})
{\"{a}}quivalent zur Erf{\"{u}}llbarkeit ist.
\qedhere
\end{proof}

\begin{satz}[Untere Schranke f{\"{u}}r ALC]
\label{satz:alc-pspace-lower}
Die Erf{\"{u}}llbarkeit von \textsf{ALC}-Konzepten ist
\textsc{PSpace}-schwer.
\end{satz}

\begin{proof}[Beweisidee]
Wir reduzieren \textsc{QBF} (quantifizierte Boolesche Formeln)
auf \textsf{ALC}-Erf{\"{u}}llbarkeit.

Eine Formel $\Phi = Q_1 x_1 \ldots Q_n x_n.\phi$
(mit $Q_i \in \{\forall, \exists\}$) wird wie folgt kodiert:

\textit{Variablen:} Jede Variable $x_i$ wird durch Konzeptname
$X_i$ kodiert (wahr) und $\neg X_i$ (falsch).

\textit{Quantoren:} Ein Existenzquantor $\exists x_i$
wird durch einen Nachfolger kodiert, der entweder
$X_i$ oder $\neg X_i$ enth{\"{a}}lt:
\[
  \exists\mathit{step}_i.(X_i \sqcup \neg X_i) \sqcap \ldots
\]
Ein Allquantor $\forall x_i$ wird durch:
\[
  \forall\mathit{step}_i.(X_i \sqcup \neg X_i) \sqcap \ldots
\]

\textit{Matrixformel:} Die propositionale Formel $\phi$
wird direkt in einen \textsf{ALC}-Konzeptausdruck {\"{u}}bersetzt.

Die Reduktion hat polynomielle Gr{\"{o}}{\ss}e und zeigt:
$\Phi$ ist wahr gdw.\ das kodierte \textsf{ALC}-Konzept
erf{\"{u}}llbar ist.

Da \textsc{QBF} \textsc{PSpace}-vollst{\"{a}}ndig ist~\cite{Stockmeyer1976},
folgt die untere Schranke.
\qedhere
\end{proof}

\begin{korollar}[PSpace-Vollst{\"{a}}ndigkeit von ALC]
Die Erf{\"{u}}llbarkeit und die Subsumtion von \textsf{ALC}-Konzepten
(bez{\"{u}}glich allgemeiner TBoxen) sind \textsc{PSpace}-vollst{\"{a}}ndig.
\end{korollar}

\begin{proof}
Erf{\"{u}}llbarkeit ist \textsc{PSpace}-vollst{\"{a}}ndig nach
S{\"{a}}tzen~\ref{satz:alc-pspace-upper} und~\ref{satz:alc-pspace-lower}.
Subsumtion $C \sqsubseteq D$ ist {\"{a}}quivalent zur Unerf{\"{u}}llbarkeit
von $C \sqcap \neg D$ (Satz~\ref{satz:red-unerfuell}),
also ebenfalls \textsc{PSpace}-vollst{\"{a}}ndig.
\end{proof}

\begin{figure}[ht]
  \centering
  \includegraphics[width=0.82\textwidth]{plots/plot3_komplexitaet.pdf}
  \caption{Komplexit{\"{a}}tslandschaft der DL-Familien. Die logarithmische
           Skala veranschaulicht den dramatischen Unterschied zwischen
           den Familien. Beschriftungen geben die Entscheidungskomplexit{\"{a}}t
           der Subsumtion an.}
  \label{fig:komplexitaet}
\end{figure}

% =======================================================================
\section{Tableau-basierte Entscheidungsverfahren}
\label{sec:tableau}
% =======================================================================

\subsection{Formale Definition des Tableau-Verfahrens}

\begin{definition}[Tableau-System]
\label{def:tableau}
Ein \emph{Tableau} f{\"{u}}r ein \textsf{ALC}-Konzept $C_0$ bez{\"{u}}glich
TBox $\Tbox$ ist ein gerichteter, markierter Baum
$T = (V, E, \mathcal{L})$ mit:
\begin{itemize}
  \item einer endlichen Knotenmenge $V$,
  \item $E \subseteq V \times N_R \times V$
        (gerichtete, mit Rollennamen markierte Kanten),
  \item einer Knotenbeschriftung
        $\mathcal{L}: V \to 2^{\mathrm{sub}^+(C_0, \Tbox)}$,
        wobei $\mathrm{sub}^+(C_0, \Tbox)$ die Menge aller
        Teilkonzepte von $C_0$ und aller in $\Tbox$ vorkommenden
        Konzepte (und ihrer NNF-Negationen) ist.
\end{itemize}
Die Wurzel $v_0$ hat $\mathcal{L}(v_0) \supseteq \{C_0\}$.
\end{definition}

\begin{definition}[Clash und Vollst{\"{a}}ndigkeit]
\label{def:clash}
\begin{itemize}
  \item Ein Knoten $v$ enth{\"{a}}lt einen \emph{Clash}, wenn
        $C \in \mathcal{L}(v)$ und $\neg C \in \mathcal{L}(v)$
        f{\"{u}}r ein Konzept $C$, oder $\bot \in \mathcal{L}(v)$.
  \item Ein Tableau hei{\ss}t \emph{vollst{\"{a}}ndig}, wenn keine
        Tableau-Regel mehr anwendbar ist.
  \item Ein Tableau hei{\ss}t \emph{geschlossen}, wenn jeder
        Blattknoten einen Clash enth{\"{a}}lt.
  \item Ein Tableau hei{\ss}t \emph{offen}, wenn es vollst{\"{a}}ndig
        und nicht geschlossen ist.
\end{itemize}
\end{definition}

\begin{definition}[Tableau-Regeln]
\label{def:tableau-regeln}
Die Tableau-Regeln f{\"{u}}r \textsf{ALC} mit TBox $\Tbox$ sind:

\medskip
\begin{tabular}{@{}p{1.8cm}p{5.5cm}p{5cm}@{}}
\toprule
\textbf{Regel} & \textbf{Bedingung} & \textbf{Aktion} \\
\midrule
$\sqcap$-Regel &
  $C \sqcap D \in \mathcal{L}(v)$, $\{C,D\} \not\subseteq \mathcal{L}(v)$ &
  $\mathcal{L}(v) {:=} \mathcal{L}(v) \cup \{C,D\}$ \\[4pt]
$\sqcup$-Regel &
  $C \sqcup D \in \mathcal{L}(v)$, $C,D \notin \mathcal{L}(v)$ &
  $\mathcal{L}(v) {:=} \mathcal{L}(v) \cup \{E\}$, $E \in \{C,D\}$ \\[4pt]
$\exists$-Regel &
  $\exists r.C \in \mathcal{L}(v)$, kein $r$-Nachf.\ mit $C$ &
  Erzeuge Nachf.\ $w$: $(v,r,w)$, $C \in \mathcal{L}(w)$ \\[4pt]
$\forall$-Regel &
  $\forall r.C \in \mathcal{L}(v)$, $(v,r,w) \in E$, $C \notin \mathcal{L}(w)$ &
  $\mathcal{L}(w) {:=} \mathcal{L}(w) \cup \{C\}$ \\[4pt]
GCI-Regel &
  $D \in \mathcal{L}(v)$, $D \sqsubseteq E \in \Tbox$, $E \notin \mathcal{L}(v)$ &
  $\mathcal{L}(v) {:=} \mathcal{L}(v) \cup \{E\}$ \\
\bottomrule
\end{tabular}
\end{definition}

\begin{figure}[ht]
  \centering
  \includegraphics[width=0.92\textwidth]{plots/plot2_tableau.pdf}
  \caption{Ablauf des Tableau-Verfahrens: Ausgehend von einem Startkonzept
           werden Regeln schrittweise angewendet. Terminiert das Verfahren
           ohne Clash, ist das Konzept erf{\"{u}}llbar.}
  \label{fig:tableau}
\end{figure}

\subsection{Korrektheit und Vollst{\"{a}}ndigkeit des Tableau-Verfahrens}

\begin{satz}[Korrektheit des Tableau-Verfahrens]
\label{satz:tableau-korrektheit}
Sei $C_0$ ein \textsf{ALC}-Konzept und $\Tbox$ eine TBox.
Dann gilt:
\[
  C_0 \text{ erf{\"{u}}llbar bzgl. } \Tbox
  \iff \text{der Tableau-Algorithmus erzeugt ein offenes Tableau.}
\]
\end{satz}

\begin{proof}
\textbf{Richtung $(\Rightarrow)$ -- Vollst{\"{a}}ndigkeit:}

Angenommen, $C_0$ ist erf{\"{u}}llbar bez{\"{u}}glich $\Tbox$.
Dann existiert ein Modell $\interp{I}$ und $d_0 \in C_0^\interp{I}$
mit $\interp{I} \models \Tbox$.

Wir konstruieren ein offenes Tableau, indem wir die
Modellstruktur von $\interp{I}$ ``unrolled'':

Definiere den Tableau-Baum $T$ induktiv:
\begin{itemize}
  \item Wurzel $v_0$ mit $\mathcal{L}(v_0) = \{C \in \mathrm{sub}^+(C_0,\Tbox) \mid d_0 \in C^\interp{I}\}$.
  \item F{\"{u}}r jeden Knoten $v$ mit Modellelement $d_v$:
        Falls $\exists r.C \in \mathcal{L}(v)$,
        w{\"{a}}hle $e$ mit $(d_v, e) \in r^\interp{I}$ und $e \in C^\interp{I}$
        (existiert nach Semantik), erzeuge Nachfolger $w$ mit
        $\mathcal{L}(w) = \{D \in \mathrm{sub}^+\mid e \in D^\interp{I}\}$
        und $d_w = e$.
\end{itemize}

\textit{Claim:} Dieses Tableau ist vollst{\"{a}}ndig (alle Regeln sind
erf{\"{u}}llt) und clashfrei.

\textit{Vollst{\"{a}}ndigkeit der Regeln:}
Die $\sqcap$-Regel ist erf{\"{u}}llt: Falls $C \sqcap D \in \mathcal{L}(v)$,
dann $d_v \in (C \sqcap D)^\interp{I} = C^\interp{I} \cap D^\interp{I}$,
also $d_v \in C^\interp{I}$ und $d_v \in D^\interp{I}$, d.\,h.\
$C, D \in \mathcal{L}(v)$.

Die $\forall$-Regel ist erf{\"{u}}llt: Falls $\forall r.C \in \mathcal{L}(v)$
und $(v,r,w) \in E$, dann $(d_v, d_w) \in r^\interp{I}$.
Da $d_v \in (\forall r.C)^\interp{I}$, gilt $d_w \in C^\interp{I}$,
also $C \in \mathcal{L}(w)$.

Die GCI-Regel: Falls $D \in \mathcal{L}(v)$ und
$D \sqsubseteq E \in \Tbox$, dann $d_v \in D^\interp{I}$
und da $\interp{I} \models D \sqsubseteq E$, folgt
$d_v \in E^\interp{I}$, also $E \in \mathcal{L}(v)$.

\textit{Clashfreiheit:}
Angenommen, $C, \neg C \in \mathcal{L}(v)$.
Dann $d_v \in C^\interp{I}$ und $d_v \in (\neg C)^\interp{I} = \Dom{I} \setminus C^\interp{I}$ -- Widerspruch.
Also ist das Tableau offen.

\textbf{Richtung $(\Leftarrow)$ -- Korrektheit (Soundness):}

Sei $T$ ein vollst{\"{a}}ndiges, offenes Tableau f{\"{u}}r $C_0$.
Wir konstruieren ein Modell $\interp{I}$ direkt aus $T$:

\begin{itemize}
  \item $\Dom{I} = V$ (Knoten als Dom{\"{a}}nenelemente).
  \item $A^\interp{I} = \{v \in V \mid A \in \mathcal{L}(v)\}$
        f{\"{u}}r alle $A \in N_C$.
  \item $r^\interp{I} = \{(v,w) \mid (v,r,w) \in E\}$
        f{\"{u}}r alle $r \in N_R$.
  \item $a^\interp{I} = v_0$ f{\"{u}}r alle Individuennamen $a$.
\end{itemize}

\textit{Claim:} $\interp{I} \models \Tbox$ und
$v_0 \in C_0^\interp{I}$.

\textit{Claim: F{\"{u}}r alle Knoten $v$ und alle
$D \in \mathcal{L}(v)$ gilt $v \in D^\interp{I}$.}

Beweis durch Induktion {\"{u}}ber den Aufbau von $D$:

\textit{Basisfall $D = A$:} $v \in A^\interp{I}$ per Definition. \checkmark

\textit{$D = \top$:} $v \in \Dom{I}$. \checkmark

\textit{$D = \neg A$:} Da das Tableau clashfrei ist,
gilt $A \notin \mathcal{L}(v)$, also $v \notin A^\interp{I}$,
d.\,h.\ $v \in (\neg A)^\interp{I}$. \checkmark

\textit{$D = C_1 \sqcap C_2$:}
Da das Tableau vollst{\"{a}}ndig ist ($\sqcap$-Regel), gilt
$C_1, C_2 \in \mathcal{L}(v)$.
Per IH $v \in C_1^\interp{I}$ und $v \in C_2^\interp{I}$,
also $v \in (C_1 \sqcap C_2)^\interp{I}$. \checkmark

\textit{$D = C_1 \sqcup C_2$:}
Da vollst{\"{a}}ndig ($\sqcup$-Regel): $C_1 \in \mathcal{L}(v)$
oder $C_2 \in \mathcal{L}(v)$.
Per IH $v \in C_1^\interp{I}$ oder $v \in C_2^\interp{I}$,
also $v \in (C_1 \sqcup C_2)^\interp{I}$. \checkmark

\textit{$D = \exists r.C$:}
Da vollst{\"{a}}ndig ($\exists$-Regel): Es gibt Nachfolger $w$
mit $(v,r,w) \in E$ und $C \in \mathcal{L}(w)$.
Per IH $w \in C^\interp{I}$ und $(v,w) \in r^\interp{I}$,
also $v \in (\exists r.C)^\interp{I}$. \checkmark

\textit{$D = \forall r.C$:}
F{\"{u}}r jeden $r$-Nachfolger $w$ von $v$:
Da vollst{\"{a}}ndig ($\forall$-Regel): $C \in \mathcal{L}(w)$.
Per IH $w \in C^\interp{I}$, also $v \in (\forall r.C)^\interp{I}$. \checkmark

Da $C_0 \in \mathcal{L}(v_0)$, gilt $v_0 \in C_0^\interp{I}$.
F{\"{u}}r die TBox: Die GCI-Regel garantiert, dass alle GCI-Konsequenzen
in den Beschriftungen eingetragen sind, also $\interp{I} \models \Tbox$.
\qedhere
\end{proof}

\begin{satz}[Terminierung des Tableau-Algorithmus]
\label{satz:terminierung}
Der Tableau-Algorithmus mit Pair-wise Blocking terminiert
f{\"{u}}r jede Eingabe $(C_0, \Tbox)$.
\end{satz}

\begin{proof}
Definiere die \emph{Tiefe} eines Tableau-Knotens $v$ als die
Lange des Pfades von der Wurzel zu $v$.

\textbf{Blocking:} Ein Knoten $v$ wird \emph{geblockt} durch
einen Vorfahren $u$, wenn $\mathcal{L}(u) = \mathcal{L}(v)$.
Geblockte Knoten erzeugen keine Nachfolger.

\textbf{Endlichkeit der Beschriftungen:}
Jede Beschriftung $\mathcal{L}(v) \subseteq \mathrm{sub}^+(C_0, \Tbox)$.
Da $\mathrm{sub}^+(C_0, \Tbox)$ endlich ist (mit
$|\mathrm{sub}^+(C_0, \Tbox)| \leq 2(|C_0| + |\Tbox|)$),
gibt es h{\"{o}}chstens $2^{2(|C_0|+|\Tbox|)}$ verschiedene Beschriftungen.

\textbf{Eindeutigkeit der Beschriftungen auf einem Pfad:}
Auf jedem Pfad von der Wurzel zu einem Blatt sind alle
Knotenbeschriftungen paarweise verschieden:
Denn die Existenz-Regel erzeugt neue Beschriftungen, die mindestens
ein neues Konzept enthalten. Gibt es denselben Beschriftungswert
zweimal auf einem Pfad, wird der tiefere Knoten geblockt.

\textbf{Tiefenschranke:} Daher hat jeder Pfad L{\"{a}}nge h{\"{o}}chstens
$2^{2(|C_0|+|\Tbox|)}$.

\textbf{Verzweigungsgrad:} An jedem Knoten k{\"{o}}nnen h{\"{o}}chstens
$|\mathrm{sub}^+(C_0,\Tbox)|$ Nachfolger erzeugt werden.

\textbf{Endlichkeit des Baums:} Der Baum hat endliche Tiefe
und endlichen Verzweigungsgrad, also endlich viele Knoten.
Jede Regelanwendung vergr{\"{o}}{\ss}ert entweder die Beschriftung
eines Knotens (endlich oft m{\"{o}}glich) oder erzeugt einen neuen Knoten
(endlich viele). Also terminiert der Algorithmus.
\qedhere
\end{proof}

\begin{beispiel}[Tableau-Konstruktion]
Sei $C_0 = \exists r.A \sqcap \forall r.\neg A$ (sollte unerf{\"{u}}llbar sein).
Tableau-Aufbau:
\begin{enumerate}
  \item Wurzel $v_0$: $\mathcal{L}(v_0) = \{\exists r.A, \forall r.\neg A\}$.
  \item $\exists$-Regel: Erzeuge $v_1$ mit $(v_0, r, v_1)$ und
        $\mathcal{L}(v_1) = \{A\}$.
  \item $\forall$-Regel: Da $(v_0,r,v_1)$ und $\forall r.\neg A \in \mathcal{L}(v_0)$:
        $\neg A \in \mathcal{L}(v_1)$.
  \item $\mathcal{L}(v_1) = \{A, \neg A\}$ -- \textbf{Clash!}
\end{enumerate}
Alle {\"{A}}ste f{\"{u}}hren zu einem Clash; das Tableau ist geschlossen.
Also ist $C_0$ unerf{\"{u}}llbar. \checkmark
\end{beispiel}

% =======================================================================
\section{Konzepthierarchie, Subsumtion und Instanzpr{\"u}fung}
\label{sec:subsumtion}
% =======================================================================

\subsection{Formale Eigenschaften der Subsumtionsrelation}

\begin{definition}[Subsumtionsrelation]
$C$ wird von $D$ \emph{subsumiert} bez{\"{u}}glich $\Tbox$
(geschrieben $C \sqsubseteq_\Tbox D$), wenn f{\"{u}}r alle
Interpretationen $\interp{I}$ mit $\interp{I} \models \Tbox$:
$C^\interp{I} \subseteq D^\interp{I}$.
\end{definition}

\begin{satz}[Pr{\"{a}}ordnungseigenschaften]
\label{satz:praeordnung}
Die Subsumtionsrelation $\sqsubseteq_\Tbox$ ist eine Pr{\"{a}}ordnung
auf $\mathcal{C}_{\textsf{ALC}}$:
\begin{enumerate}[label=(\roman*)]
  \item \textbf{Reflexivit{\"{a}}t:} $C \sqsubseteq_\Tbox C$ f{\"{u}}r alle $C$.
  \item \textbf{Transitivit{\"{a}}t:} Wenn $C \sqsubseteq_\Tbox D$
        und $D \sqsubseteq_\Tbox E$, dann $C \sqsubseteq_\Tbox E$.
\end{enumerate}
Ferner gilt:
\begin{enumerate}[label=(\roman*), resume]
  \item \textbf{Monotonie der Operatoren:} Wenn $C \sqsubseteq_\Tbox D$,
        dann $C \sqcap E \sqsubseteq_\Tbox D \sqcap E$ f{\"{u}}r alle $E$.
  \item \textbf{Subsumtion und Konjunktion:}
        $C \sqcap D \sqsubseteq_\Tbox C$ und
        $C \sqcap D \sqsubseteq_\Tbox D$.
  \item \textbf{Subsumtion und Disjunktion:}
        $C \sqsubseteq_\Tbox C \sqcup D$ und
        $D \sqsubseteq_\Tbox C \sqcup D$.
\end{enumerate}
\end{satz}

\begin{proof}
Alle Aussagen folgen direkt aus der mengentheoretischen Semantik.

\textbf{(i) Reflexivit{\"{a}}t:}
F{\"{u}}r jede Interpretation $\interp{I}$ gilt $C^\interp{I} \subseteq C^\interp{I}$
(Reflexivit{\"{a}}t der Mengenininklusion).

\textbf{(ii) Transitivit{\"{a}}t:}
Seien $C \sqsubseteq_\Tbox D$ und $D \sqsubseteq_\Tbox E$.
F{\"{u}}r jedes Modell $\interp{I}$:
$C^\interp{I} \subseteq D^\interp{I} \subseteq E^\interp{I}$.
Transitivit{\"{a}}t der Mengenininklusion liefert $C^\interp{I} \subseteq E^\interp{I}$.

\textbf{(iii) Monotonie:}
Aus $C^\interp{I} \subseteq D^\interp{I}$ folgt
$(C \sqcap E)^\interp{I} = C^\interp{I} \cap E^\interp{I}
\subseteq D^\interp{I} \cap E^\interp{I}
= (D \sqcap E)^\interp{I}$.

\textbf{(iv):} $(C \sqcap D)^\interp{I} = C^\interp{I} \cap D^\interp{I} \subseteq C^\interp{I}$.

\textbf{(v):} $C^\interp{I} \subseteq C^\interp{I} \cup D^\interp{I} = (C \sqcup D)^\interp{I}$.
\qedhere
\end{proof}

\begin{proposition}[Subsumtion durch Negation]
F{\"{u}}r \textsf{ALC}-Konzepte gilt:
$C \sqsubseteq_\Tbox D$ genau dann, wenn
$\top \sqsubseteq_\Tbox \neg C \sqcup D$.
\end{proposition}

\begin{proof}
$C \sqsubseteq_\Tbox D$
iff f{\"{u}}r alle Modelle $\interp{I}$: $C^\interp{I} \subseteq D^\interp{I}$
iff $C^\interp{I} \cap (\Dom{I} \setminus D^\interp{I}) = \emptyset$
iff $\Dom{I} \subseteq (\Dom{I} \setminus C^\interp{I}) \cup D^\interp{I}
= (\neg C \sqcup D)^\interp{I}$
iff $\top \sqsubseteq_\Tbox \neg C \sqcup D$.
\end{proof}

\subsection{Instanz{\"{u}}berpr{\"{u}}fung und Realisierung}

\begin{definition}[Instanz{\"{u}}berpr{\"{u}}fung]
Gegeben $\KB = (\Tbox, \Abox)$, $a \in N_I$ und Konzept $C$.
$a$ ist eine \emph{Instanz} von $C$ bez{\"{u}}glich $\KB$
(geschrieben $\KB \models C(a)$), wenn f{\"{u}}r alle Modelle $\interp{I}$
von $\KB$: $a^\interp{I} \in C^\interp{I}$.
\end{definition}

\begin{satz}[Instanzpr{\"{u}}fung durch Unerf{\"{u}}llbarkeit]
\label{satz:instanz}
$\KB \models C(a)$ gdw.\ $(\Tbox, \Abox \cup \{\neg C(a)\})$
inkonsistent ist.
\end{satz}

\begin{proof}
Dies ist Satz~\ref{satz:red-unerfuell}(ii), dort vollst{\"{a}}ndig bewiesen.
\end{proof}

\begin{satz}[Komplexit{\"{a}}t der Instanzpr{\"{u}}fung]
Die Instanzpr{\"{u}}fung in \textsf{ALC} ist \textsc{PSpace}-vollst{\"{a}}ndig
(kombinierte Komplexit{\"{a}}t) und \textsc{P}-vollst{\"{a}}ndig
(Datenkomplexit{\"{a}}t, wenn $\Tbox$ fest).
\end{satz}

\begin{proof}[Beweisidee]
Obere Schranke: Instanzpr{\"{u}}fung reduziert auf Unerf{\"{u}}llbarkeit
(Satz~\ref{satz:instanz}), die in \textsc{PSpace} liegt.

Datenkomplexit{\"{a}}t: Wenn $\Tbox$ fest ist, kann man die Anfrage
durch ein polynomielles Datenbankproblem kodieren (Rewriting-Technik).
\end{proof}

% =======================================================================
\section{Erweiterte DL-Familien: SHIQ und SROIQ}
\label{sec:shiq-sroiq}
% =======================================================================

\subsection{SHIQ: Transitive Rollen, Rollenhierarchien und Z{\"{a}}hlen}

\begin{definition}[SHIQ-Syntax]
\textsf{SHIQ} erweitert \textsf{ALC} um:
\begin{enumerate}[label=(\roman*)]
  \item \textbf{Transitive Rollen (S):}
        Rollenaxiome $\mathrm{Trans}(r)$;
        Semantik: $r^\interp{I}$ ist transitiv.
  \item \textbf{Rollenhierarchien (H):}
        GCIs $r \sqsubseteq s$ f{\"{u}}r Rollennamen $r,s$;
        Semantik: $r^\interp{I} \subseteq s^\interp{I}$.
  \item \textbf{Inverse Rollen (I):}
        $r^-$; Semantik: $(r^-)^\interp{I} = \{(e,d) \mid (d,e) \in r^\interp{I}\}$.
  \item \textbf{Qualifizierte Z{\"{a}}hlrestriktionen (Q):}
        $\geq n\,r.C$ und $\leq n\,r.C$;
        Semantik: $(\geq n\,r.C)^\interp{I} = \{d \mid |\{e \mid (d,e) \in r^\interp{I}, e \in C^\interp{I}\}| \geq n\}$.
\end{enumerate}
\end{definition}

\begin{satz}[Korrektheit der Z{\"{a}}hlrestriktions-Semantik]
\label{satz:zaehlung}
F{\"{u}}r jede Interpretation $\interp{I}$ gilt:
\[
  (\leq n\,r.C)^\interp{I} =
  \Dom{I} \setminus (\geq n+1\,r.C)^\interp{I}.
\]
\end{satz}

\begin{proof}
Per Definition:
\begin{align*}
  (\leq n\,r.C)^\interp{I}
  &= \{d \mid |\{e \mid (d,e) \in r^\interp{I}, e \in C^\interp{I}\}| \leq n\} \\
  &= \Dom{I} \setminus \{d \mid |\{e \mid (d,e) \in r^\interp{I},
    e \in C^\interp{I}\}| \geq n+1\} \\
  &= \Dom{I} \setminus (\geq n+1\,r.C)^\interp{I}. \qedhere
\end{align*}
\end{proof}

\begin{satz}[ExpTime-Vollst{\"{a}}ndigkeit von SHIQ, \cite{Horrocks2000}]
Die Erf{\"{u}}llbarkeit in \textsf{SHIQ} bez{\"{u}}glich allgemeiner TBoxen ist
\textsc{ExpTime}-vollst{\"{a}}ndig.
\end{satz}

\begin{proof}[Beweisidee]
\textbf{Obere Schranke:}
Das Tableau-Verfahren f{\"{u}}r \textsf{SHIQ} erzeugt B{\"{a}}ume der
exponentiellen Tiefe. Transitive Rollen erfordern eine
spezielle Form des Blocking (``lazy unfolding''),
um korrekte Terminierung zu garantieren.
Ein deterministischer Algorithmus verarbeitet den exponentiell
gro{\ss}en Baum in exponentieller Zeit -- die Korrektheit des
Verfahrens zeigt $\in \textsc{ExpTime}$.

\textbf{Untere Schranke:}
Durch Reduktion von \textsc{ALC}-TBox-Erf{\"{u}}llbarkeit
(die bereits \textsc{ExpTime}-schwer ist f{\"{u}}r allgemeine TBoxen,
vgl.~\cite{Baader2007})
auf \textsf{SHIQ}-Erf{\"{u}}llbarkeit.
\end{proof}

\subsection{SROIQ und OWL 2}

\begin{definition}[SROIQ]
\textsf{SROIQ} erweitert \textsf{SHIQ} um:
\begin{enumerate}[label=(\roman*)]
  \item \textbf{Nominale/OneOf (O):}
        $\{a_1, \ldots, a_k\}$ mit $a_i \in N_I$.
        Semantik: $\{a_1,\ldots,a_k\}^\interp{I} = \{a_1^\interp{I},\ldots,a_k^\interp{I}\}$.
  \item \textbf{Komplexe Rollenaxiome (RBox):}
        $r_1 \circ \ldots \circ r_k \sqsubseteq s$ mit $r_i, s \in N_R$.
        Semantik: $(r_1 \circ \cdots \circ r_k)^\interp{I}
        \subseteq s^\interp{I}$.
  \item \textbf{Disjunkte Rollen:} $\mathrm{Dis}(r,s)$,
        Semantik: $r^\interp{I} \cap s^\interp{I} = \emptyset$.
\end{enumerate}
\end{definition}

\begin{satz}[NExpTime-Vollst{\"{a}}ndigkeit, \cite{Horrocks2006}]
Die Erf{\"{u}}llbarkeit in \textsf{SROIQ} ist \textsc{NExpTime}-vollst{\"{a}}ndig.
\end{satz}

Die vollst{\"{a}}ndige Beweisf{\"{u}}hrung erfordert die Analyse von
hypertableaux-Regeln f{\"{u}}r komplexe Rollenaxiome sowie
den Nachweis, dass gewisse Einschr{\"{a}}nkungen der RBox
(``regular'' role hierarchies) notwendig f{\"{u}}r Entscheidbarkeit sind.
Details finden sich in~\cite{Horrocks2006}.

% =======================================================================
\section{Lernen von Beschreibungslogik-Konzepten}
\label{sec:lernen}
% =======================================================================

\subsection{Das Lernproblem}

\begin{definition}[Lernproblem f{\"{u}}r DL]
\label{def:lernproblem}
Gegeben:
\begin{itemize}
  \item eine Wissensbasis $\KB = (\Tbox, \Abox)$,
  \item eine Menge positiver Beispiele $E^+ \subseteq N_I$,
  \item eine Menge negativer Beispiele $E^- \subseteq N_I$,
    mit $E^+ \cap E^- = \emptyset$.
\end{itemize}
Gesucht: ein Konzept $C$, das die Beispiele \emph{trennt}:
\begin{align*}
  \KB &\models C(a) \text{ f{\"{u}}r alle } a \in E^+, \\
  \KB &\not\models C(a) \text{ f{\"{u}}r alle } a \in E^-.
\end{align*}
\end{definition}

\subsection{Most Specific Concept (MSC)}

\begin{definition}[Most Specific Concept]
\label{def:msc}
Ein \textsf{EL}-Konzept $C$ ist das \emph{most specific concept}
(MSC) eines Individuums $a$ bez{\"{u}}glich $\KB$, wenn:
\begin{enumerate}[label=(\roman*)]
  \item $\KB \models C(a)$ (Korrektheit),
  \item F{\"{u}}r alle \textsf{EL}-Konzepte $D$:
        Wenn $\KB \models D(a)$, dann $\Tbox \models C \sqsubseteq D$
        (Spezifit{\"{a}}t).
\end{enumerate}
\end{definition}

\begin{satz}[Existenz und Berechenbarkeit des MSC]
\label{satz:msc-exist}
In \textsf{EL} existiert das MSC jedes Individuums $a$
bez{\"{u}}glich jeder konsistenten Wissensbasis $\KB = (\Tbox, \Abox)$
und ist in Polynomialzeit berechenbar.
\end{satz}

\begin{proof}
\textbf{Konstruktion:}
Wir konstruieren $\mathrm{MSC}(a)$ durch Verfolgung der
ABox-Struktur:

Definiere f{\"{u}}r jedes Individuum $b \in N_I$ induktiv
den Konzeptausdruck $T(b)$:
\begin{align*}
  T(b) &= \bigcap_{A(b) \in \Abox} A
          \;\sqcap\;
          \bigcap_{r(b,c) \in \Abox} \exists r.T(c).
\end{align*}

Diese Konstruktion ist wohldefiniert, wenn die ABox azyklisch ist.
Bei zyklischen ABoxen muss der kleinste Fixpunkt gebildet werden.

\textbf{Korrektheit:}
In jedem Modell $\interp{I}$ von $\KB$ gilt:
\begin{itemize}
  \item F{\"{u}}r jeden Konzeptnamen $A$ mit $A(a) \in \Abox$:
        $a^\interp{I} \in A^\interp{I}$, also
        $a^\interp{I} \in (\bigcap_{A(a) \in \Abox} A)^\interp{I}$.
  \item F{\"{u}}r jede Rollenassertion $r(a,b) \in \Abox$:
        $(a^\interp{I}, b^\interp{I}) \in r^\interp{I}$
        und $b^\interp{I} \in T(b)^\interp{I}$ (per IH),
        also $a^\interp{I} \in (\exists r.T(b))^\interp{I}$.
\end{itemize}
Damit $\KB \models T(a)(a)$, d.\,h.\ $\mathrm{MSC}(a) = T(a)$
ist korrekt.

\textbf{Spezifit{\"{a}}t:}
Sei $D$ ein \textsf{EL}-Konzept mit $\KB \models D(a)$.
Wir zeigen $\Tbox \models T(a) \sqsubseteq D$ durch
Strukturinduktion:

Falls $D = A$: Dann muss $A(a) \in \Abox$ gelten
(oder durch Schlussfolgern aus der TBox folgen).
In diesem Fall enth{\"{a}}lt $T(a)$ den Konjunkt $A$,
also $T(a) \sqsubseteq A$ trivial.

Falls $D = D_1 \sqcap D_2$:
$\KB \models D_1(a)$ und $\KB \models D_2(a)$.
Per IH $\Tbox \models T(a) \sqsubseteq D_1$
und $\Tbox \models T(a) \sqsubseteq D_2$,
also $\Tbox \models T(a) \sqsubseteq D_1 \sqcap D_2$.

Falls $D = \exists r.E$:
$\KB \models \exists r.E(a)$, also gibt es $b$
mit $r(a,b) \in \Abox$ und $\KB \models E(b)$.
Per IH $\Tbox \models T(b) \sqsubseteq E$.
Da $\exists r.T(b)$ ein Konjunkt von $T(a)$ ist,
gilt $\Tbox \models T(a) \sqsubseteq \exists r.T(b) \sqsubseteq \exists r.E$.

\textbf{Polynomialzeit:}
Die Konstruktion durchl{\"{a}}uft die ABox-Struktur und hat
Laufzeit $O(|\Abox|)$.
\qedhere
\end{proof}

\subsection{Least Common Subsumer (LCS)}

\begin{definition}[Least Common Subsumer]
\label{def:lcs}
Das \emph{least common subsumer} (LCS) zweier \textsf{EL}-Konzepte
$C$ und $D$ bez{\"{u}}glich $\Tbox$ ist ein \textsf{EL}-Konzept
$\mathrm{lcs}(C, D)$ mit:
\begin{enumerate}[label=(\roman*)]
  \item $C \sqsubseteq_\Tbox \mathrm{lcs}(C,D)$
        und $D \sqsubseteq_\Tbox \mathrm{lcs}(C,D)$
        (gemeinsamer Oberbegriff),
  \item F{\"{u}}r alle \textsf{EL}-Konzepte $E$:
        Falls $C \sqsubseteq_\Tbox E$ und $D \sqsubseteq_\Tbox E$,
        dann $\mathrm{lcs}(C,D) \sqsubseteq_\Tbox E$
        (Minimalit{\"{a}}t).
\end{enumerate}
\end{definition}

\begin{satz}[Existenz und Berechnung des LCS in EL, \cite{Baader1999}]
\label{satz:lcs}
In \textsf{EL} existiert f{\"{u}}r je zwei Konzepte $C$ und $D$
ein LCS, der durch Produktkonstruktion berechnet werden kann
und polynomielle Gr{\"{o}}{\ss}e in $|C| \cdot |D|$ hat.
\end{satz}

\begin{proof}
\textbf{Schritt 1 -- Beschreibungsbaumdarstellung:}
Jedes \textsf{EL}-Konzept $C$ kann als gerichteter,
geordneter, mit Konzeptnamen beschrifteter Baum $\mathcal{T}_C$
dargestellt werden (``description tree''):
\begin{itemize}
  \item F{\"{u}}r $C = A$: einkn{\"{o}}tiger Baum mit Beschriftung $\{A\}$.
  \item F{\"{u}}r $C = C_1 \sqcap C_2$:
        Wurzel mit Beschriftung $\mathrm{labels}(C_1) \cup \mathrm{labels}(C_2)$
        und den Nachfolgern von $\mathcal{T}_{C_1}$ und $\mathcal{T}_{C_2}$.
  \item F{\"{u}}r $C = \exists r.D$:
        Wurzel mit $r$-markierter Kante zu $\mathcal{T}_D$.
\end{itemize}

\textbf{Schritt 2 -- Produktkonstruktion:}
Der Produktgraph $\mathcal{T}_C \times \mathcal{T}_D$
ist definiert als:
\begin{itemize}
  \item Knotenmenge: $V(\mathcal{T}_C) \times V(\mathcal{T}_D)$,
  \item Kante $((u,v), r, (u',v'))$ existiert, wenn
        $(u,r,u') \in \mathcal{T}_C$ und $(v,r,v') \in \mathcal{T}_D$,
  \item Knotenbeschriftung:
        $\mathrm{label}(u,v) = \mathrm{label}(u) \cap \mathrm{label}(v)$.
\end{itemize}
Die Wurzel ist $(r_C, r_D)$ mit den Wurzeln $r_C$ von $\mathcal{T}_C$
und $r_D$ von $\mathcal{T}_D$.

\textbf{Schritt 3 -- Korrektheit:}
Der Produktgraph repr{\"{a}}sentiert $\mathrm{lcs}(C,D)$.

\textit{(i) Gemeinsamer Oberbegriff:}
Es gibt einen Baum-Homomorphismus $h_C: \mathcal{T}_{lcs} \to \mathcal{T}_C$
(definiert durch $h_C(u,v) = u$) und analog $h_D$.
Ein Homomorphismus zwischen Beschreibungsb{\"{a}}umen entspricht
einer Subsumtionsbeziehung in \textsf{EL}~\cite{Baader2007}.
Also $\mathrm{lcs}(C,D) \sqsubseteq C$ und $\mathrm{lcs}(C,D) \sqsubseteq D$.

\textit{(ii) Minimalit{\"{a}}t:}
Sei $E$ ein \textsf{EL}-Konzept mit $C \sqsubseteq E$ und $D \sqsubseteq E$.
Dann gibt es Homomorphismen $h: \mathcal{T}_E \to \mathcal{T}_C$
und $g: \mathcal{T}_E \to \mathcal{T}_D$.
Der Produkthomomorphismus $(h,g): \mathcal{T}_E \to \mathcal{T}_C \times \mathcal{T}_D$
zeigt $\mathrm{lcs}(C,D) \sqsubseteq E$.

\textbf{Schritt 4 -- Gr{\"{o}}{\ss}e:}
$|V(\mathcal{T}_C \times \mathcal{T}_D)|
\leq |V(\mathcal{T}_C)| \cdot |V(\mathcal{T}_D)|
\leq |C| \cdot |D|$.
\qedhere
\end{proof}

\begin{beispiel}[LCS-Berechnung]
Seien $C = \mathit{Patient} \sqcap \exists\mathit{hasDisease}.\mathit{Cancer}$
und $D = \mathit{Patient} \sqcap \exists\mathit{hasDisease}.\mathit{Flu}$.

Beschreibungsb{\"{a}}ume:
\[
  \mathcal{T}_C:\quad [\mathit{Patient}] \xrightarrow{\mathit{hasDisease}} [\mathit{Cancer}], \quad
  \mathcal{T}_D:\quad [\mathit{Patient}] \xrightarrow{\mathit{hasDisease}} [\mathit{Flu}].
\]

Produktkonstruktion:
Knotenmenge: $\{(r_C, r_D), (c, d)\}$
wobei $c$ der $\mathit{hasDisease}$-Nachfolger in $\mathcal{T}_C$
und $d$ der in $\mathcal{T}_D$ ist.

Knotenbeschriftungen:
$(r_C, r_D)$: $\{\mathit{Patient}\} \cap \{\mathit{Patient}\} = \{\mathit{Patient}\}$,
$(c, d)$: $\{\mathit{Cancer}\} \cap \{\mathit{Flu}\} = \emptyset$ (nur $\top$).

Ergebnis:
$\mathrm{lcs}(C,D) = \mathit{Patient} \sqcap \exists\mathit{hasDisease}.\top
\equiv \mathit{Patient} \sqcap \exists\mathit{hasDisease}.\mathit{Disease}$
(falls $\mathit{Cancer} \sqsubseteq \mathit{Disease}$ und $\mathit{Flu} \sqsubseteq \mathit{Disease}$ in $\Tbox$).
\end{beispiel}

\begin{figure}[ht]
  \centering
  \includegraphics[width=0.90\textwidth]{plots/plot5_lernen.pdf}
  \caption{Schematischer Ablauf des Konzeptlernens:
           Aus den MSCs der positiven Beispiele wird durch LCS-Bildung
           das spezifischste gemeinsame Konzept berechnet.}
  \label{fig:lernen}
\end{figure}

% =======================================================================
\section{Nicht-monotones Schlussfolgern}
\label{sec:nichtmonoton}
% =======================================================================

\subsection{Motivierende Beispiele und Probleme}

Klassische Beschreibungslogiken sind monoton:

\begin{satz}[Monotonie]
\label{satz:monotonie}
Wenn $\KB \models C(a)$ und $\KB \subseteq \KB'$
(d.\,h.\ $\Tbox \subseteq \Tbox'$ und $\Abox \subseteq \Abox'$),
dann $\KB' \models C(a)$.
\end{satz}

\begin{proof}
Jedes Modell von $\KB'$ ist auch Modell von $\KB$
(mehr Axiome = strengere Bedingungen; irrt{\"{u}}mlicherweise:
$\KB' \supseteq \KB$ bedeutet, alle Modelle von $\KB'$ erf{\"{u}}llen auch $\KB$).
Aber: $\KB \subseteq \KB'$ bedeutet, jedes Modell von $\KB'$
ist ein Modell von $\KB$, denn $\KB'$ macht zus{\"{a}}tzliche Forderungen.
Falls $\KB \models C(a)$, dann gilt in jedem Modell von $\KB$,
also in jedem Modell von $\KB'$: $a^\interp{I} \in C^\interp{I}$.
\qedhere
\end{proof}

Monotonie ist in der Praxis problematisch: Man m{\"{o}}chte
\emph{typische} Eigenschaften angeben, die durch spezifischere
Klassen {\"{u}}berschrieben werden k{\"{o}}nnen.

\subsection{Defeasible Inklusionen}

\begin{definition}[Pr{\"{a}}ferenzmodell und Typikalit{\"{a}}t]
Ein \emph{Pr{\"{a}}ferenzmodell} ist ein Tupel
$\interp{P} = (\Dom{P}, \cdot^\interp{P}, \prec)$,
wobei $\prec \subseteq \Dom{P} \times \Dom{P}$
eine strikt azyklische, wohlgegr{\"{u}}ndete Ordnung ist
(\emph{Pr{\"{a}}ferenzordnung}, ``typischer ist kleiner'').

Die Erweiterung des Typikalit{\"{a}}tsoperators $\mathbf{T}$:
\[
  \mathbf{T}(C)^\interp{P} =
  \bigl\{d \in C^\interp{P} \;\big|\;
  \nexists e \in C^\interp{P}.\, e \prec d\bigr\},
\]
d.\,h.\ die minimalen (typischsten) Elemente von $C$.
\end{definition}

\begin{satz}[Defeasible Modus Ponens]
\label{satz:defeasible-mp}
In Pr{\"{a}}ferenzmodellen gilt:
Wenn $\mathbf{T}(C) \sqsubseteq D$ und $a$ ist
ein typisches $C$-Element (d.\,h.\ $a \in \mathbf{T}(C)^\interp{P}$),
dann $a \in D^\interp{P}$.

Dies gilt jedoch \emph{nicht global} f{\"{u}}r alle Instanzen von $C$.
\end{satz}

\begin{proof}
Direkt aus der Definition: $\mathbf{T}(C)^\interp{P} \subseteq D^\interp{P}$
(da $\mathbf{T}(C) \sqsubseteq D$) und $a \in \mathbf{T}(C)^\interp{P}$
impliziert $a \in D^\interp{P}$.

Das ``nicht global'': Ein nichtminimales Element $b \in C^\interp{P}$
mit $b \notin \mathbf{T}(C)^\interp{P}$ muss nicht in $D^\interp{P}$
liegen.
\end{proof}

\begin{beispiel}[Ausnahmen durch Subklassen]
\begin{align*}
  \mathbf{T}(\mathit{Bird}) &\sqsubseteq \mathit{CanFly}, \\
  \mathit{Penguin} &\sqsubseteq \mathit{Bird}, \\
  \mathbf{T}(\mathit{Penguin}) &\sqsubseteq \neg\mathit{CanFly}.
\end{align*}
Ein typischer Vogel kann fliegen; ein typischer Pinguin nicht.
Da Pinguine im Pr{\"{a}}ferenzmodell ``spezifischer'' (kleiner in $\prec$)
als gew{\"{o}}hnliche V{\"{o}}gel sind, {\"{u}}berschreibt die Pinguin-Regel die Vogel-Regel.
\end{beispiel}

% =======================================================================
\section{Relaxierte Anfrageauswertung}
\label{sec:anfragen}
% =======================================================================

\subsection{Konjunktive Anfragen}

\begin{definition}[Konjunktive Anfrage (CQ)]
Eine \emph{konjunktive Anfrage} $q(\vec{x})$ {\"{u}}ber einer Signatur $\Sigma$
ist ein Ausdruck der Form:
\[
  q(\vec{x}) = \exists \vec{y}.\,
  \bigwedge_{i=1}^{k} A_i(t_i) \wedge \bigwedge_{j=1}^{m} r_j(s_j, s_j'),
\]
wobei $\vec{x}$ die freien Variablen (\emph{Antwortterme}),
$\vec{y}$ existenziell quantifizierte Variablen sind,
und $t_i, s_j, s_j'$ Terme aus $\vec{x} \cup \vec{y} \cup N_I$.
\end{definition}

\begin{definition}[Gewisse Antwortmenge]
\label{def:cert}
Die \emph{gewisse Antwortmenge} einer CQ $q(\vec{x})$
bez{\"{u}}glich $\KB = (\Tbox, \Abox)$ ist:
\[
  \mathrm{cert}(q, \KB) =
  \bigl\{\vec{a} \in N_I^{|\vec{x}|} \;\big|\;
  \text{f{\"{u}}r jedes Modell } \interp{I} \text{ von } \KB:\,
  \interp{I} \models q(\vec{a})\bigr\}.
\]
\end{definition}

\begin{satz}[Datenkomplexit{\"{a}}t der CQ-Beantwortung in EL]
\label{satz:cq-data}
Die Beantwortung konjunktiver Anfragen bez{\"{u}}glich \textsf{EL}-TBoxen
ist in \textsc{P} (polynomiell) in der Datenkomplexit{\"{a}}t.
\end{satz}

\begin{proof}[Beweisidee]
Bei fester TBox $\Tbox$ kann die Anfragebeantwortung durch
\emph{Query Rewriting} gel{\"{o}}st werden:
Die CQ $q$ wird zusammen mit $\Tbox$ in eine neue CQ $q'$
(ohne Ontologie) umgeschrieben, so dass
$\mathrm{cert}(q, (\Tbox,\Abox)) = \mathrm{cert}(q', (\emptyset, \Abox))$.

Das Rewriting $q'$ kann durch den \textsf{EL}-Subsumtionsalgorithmus
(Satz~\ref{satz:el-poly}) in polynomieller Zeit berechnet werden.
Die Auswertung von $q'$ {\"{u}}ber $\Abox$ ist in \textsc{AC}$^0 \subset \textsc{P}$.
\end{proof}

\subsection{Repair-basierte Inkonsistenztoleranz}

\begin{definition}[Repair]
\label{def:repair}
Sei $\KB = (\Tbox, \Abox)$ eine inkonsistente Wissensbasis.
Eine \emph{Repair} $\mathcal{R}$ von $\KB$ bez{\"{u}}glich $\Tbox$
ist eine maximale konsistente Teilmenge von $\Abox$:
\begin{enumerate}[label=(\roman*)]
  \item $(\Tbox, \mathcal{R})$ ist konsistent,
  \item Es gibt keine strikt gr{\"{o}}{\ss}ere Menge $\mathcal{R}' \supsetneq \mathcal{R}$
        mit $\mathcal{R}' \subseteq \Abox$ und $(\Tbox, \mathcal{R}')$ konsistent.
\end{enumerate}
Mit $\mathrm{Rep}(\KB)$ bezeichnen wir die Menge aller Repairs.
\end{definition}

\begin{satz}[Eigenschaften von Repairs]
\label{satz:repairs}
F{\"{u}}r jede inkonsistente Wissensbasis $\KB = (\Tbox, \Abox)$ gilt:
\begin{enumerate}[label=(\roman*)]
  \item $\mathrm{Rep}(\KB) \neq \emptyset$
        (Existenz mindestens einer Repair).
  \item Jede Repair $\mathcal{R}$ ist maximal konsistent.
  \item $\bigcap_{\mathcal{R} \in \mathrm{Rep}(\KB)} \mathcal{R}$
        ist der maximale konsistente Kern.
\end{enumerate}
\end{satz}

\begin{proof}
\textbf{(i) Existenz:}
$\emptyset \subseteq \Abox$ und $(\Tbox, \emptyset)$ ist konsistent
(leere ABox hat immer ein Modell: Einelementinterpretation).
Also gibt es mindestens eine konsistente Teilmenge.
Nach dem Lemma von Zorn (oder direkter endlicher Vergr{\"{o}}{\ss}erung)
gibt es eine maximale -- die Repair.

\textbf{(ii)} folgt direkt aus der Definition.

\textbf{(iii):} Der Schnitt aller Repairs enth{\"{a}}lt genau jene
Assertionen, die in jeder Repair enthalten sind -- das ist der
``gesicherte'' Teil der ABox.
\qedhere
\end{proof}

\begin{definition}[Brave und Cautious Semantik]
F{\"{u}}r eine inkonsistente Wissensbasis $\KB$ und Anfrage $q(\vec{x})$:
\begin{align*}
  \mathrm{brave}(q, \KB) &=
    \bigcup_{\mathcal{R} \in \mathrm{Rep}(\KB)} \mathrm{cert}(q, (\Tbox, \mathcal{R})), \\
  \mathrm{cautious}(q, \KB) &=
    \bigcap_{\mathcal{R} \in \mathrm{Rep}(\KB)} \mathrm{cert}(q, (\Tbox, \mathcal{R})).
\end{align*}
\end{definition}

\begin{satz}[Korrektheit der Repair-Semantik]
\label{satz:repair-korrektheit}
F{\"{u}}r eine konsistente Wissensbasis $\KB$ (mit $|\mathrm{Rep}(\KB)| = 1$,
da $\Abox$ selbst die einzige Repair ist) gilt:
\[
  \mathrm{brave}(q, \KB) = \mathrm{cautious}(q, \KB) = \mathrm{cert}(q, \KB).
\]
\end{satz}

\begin{proof}
Wenn $\KB$ konsistent ist, dann ist $\Abox$ selbst die einzige
maximale konsistente Teilmenge von $\Abox$, also
$\mathrm{Rep}(\KB) = \{\Abox\}$.
Dann:
\[
  \mathrm{brave}(q,\KB) = \mathrm{cert}(q, (\Tbox, \Abox)) = \mathrm{cert}(q,\KB)
\]
und ebenso f{\"{u}}r $\mathrm{cautious}$.
\qedhere
\end{proof}

\begin{figure}[ht]
  \centering
  \includegraphics[width=0.92\textwidth]{plots/plot6_relaxierte_anfragen.pdf}
  \caption{Relaxierte Anfrageauswertung: Links -- Vergleich von
           Standard- und relaxierten Antwortmengen f{\"{u}}r drei CQs.
           Rechts -- Visualisierung der Repair-Aufteilung bei inkonsistenter
           ABox.}
  \label{fig:relaxiert}
\end{figure}

\begin{satz}[Komplexit{\"{a}}t der Brave-Anfragebeantwortung]
Die brave repair-basierte Anfragebeantwortung f{\"{u}}r \textsf{ALC}-TBoxen
ist $\Sigma_2^p$-vollst{\"{a}}ndig in der Datenkomplexit{\"{a}}t~\cite{Bienvenu2014}.
\end{satz}

\begin{proof}[Beweisidee]
\textbf{Obere Schranke ($\Sigma_2^p$):}
Brave Semantik erfordert: ``Es gibt eine Repair $\mathcal{R}$,
sodass $\vec{a} \in \mathrm{cert}(q, (\Tbox,\mathcal{R}))$''.
Dies ist ein $\exists$-Statement {\"{u}}ber Repairs (ein NP-Orakel
w{\"{a}}hlt $\mathcal{R}$, ein co-NP-Orakel verifiziert Maximalit{\"{a}}t).
Also liegt das Problem in $\Sigma_2^p = \mathrm{NP}^{\mathrm{NP}}$.

\textbf{Untere Schranke:}
Durch Reduktion von $\Sigma_2^p$-vollst{\"{a}}ndigen Problemen
(z.\,B.\ $\exists$-QBF mit einer Ebene).
\end{proof}

% =======================================================================
\section{Beschreibungslogik und der Subgraph-Algo\-rith\-mus}
\label{sec:dl-subgraph}
% =======================================================================

\subsection{Motivation: Zwei Welten, eine Struktur}

Auf den ersten Blick scheinen Beschreibungslogiken und der
Subgraph-Algorithmus~\cite{Epp2026} v{\"o}llig verschiedene Gebiete
zu adressieren: DL ist eine formale Wissensrepresentationssprache,
der Subgraph-Algorithmus ein effizienter Graph-Vergleichsalgorithmus.
Bei n{\"a}herer Betrachtung offenbaren sich jedoch vier tiefe,
bisher nicht formal ausgearbeitete Verbindungen:

\begin{enumerate}[label=(\Roman*)]
  \item Eine \emph{DL-Interpretation} ist strukturell ein
        gerichteter, beschrifteter Graph -- und Konzept-Instanzierung
        entspricht einer \emph{Subgraph-Einbettung}.
  \item Die \emph{Signatur-Funktion} des Subgraph-Algorithmus
        kodiert exakt die Rollennachbarschaft eines ABox-Individuums
        und liefert damit einen algorithmisch auswertbaren
        Fingerabdruck f{\"u}r MSC und LCS.
  \item Das \emph{Blocking} im Tableau-Verfahren entspricht
        formal der \emph{Rotations{\"a}quivalenz} von Signatur-Vektoren.
  \item Die Beantwortung \emph{konjunktiver Anfragen} {\"u}ber einer
        ABox ist strukturell {\"a}quivalent zum
        \emph{Subgraph-Matching-Problem}.
\end{enumerate}

Alle vier Verbindungen werden in den folgenden Unterabschnitten
vollst{\"a}ndig und formal entwickelt und bewiesen.

\subsection{DL-Interpretationen als gerichtete Graphen}
\label{sec:dl-graph}

\begin{definition}[Rollengraph einer Interpretation]
\label{def:rollengraph}
Sei $\interp{I} = (\Dom{I}, \cdot^\interp{I})$ eine Interpretation
{\"u}ber Signatur $\Sigma = (N_C, N_R, N_I)$.
Der \emph{Rollengraph} $G_\interp{I} = (V, E, \lambda)$ ist definiert durch:
\begin{itemize}
  \item Knotenmenge $V = \Dom{I}$,
  \item Kantenmenge $E = \bigcup_{r \in N_R} r^\interp{I}
        \subseteq \Dom{I} \times \Dom{I}$,
  \item Knotenbeschriftung $\lambda: V \to 2^{N_C}$
        mit $\lambda(d) = \{A \in N_C \mid d \in A^\interp{I}\}$.
\end{itemize}
\end{definition}

\begin{satz}[Konzept-Instanz als Subgraph-Einbettung]
\label{satz:instanz-subgraph}
Sei $C$ ein \textsf{EL}-Konzept,
$\interp{I}$ eine Interpretation und $d \in \Dom{I}$.
Dann gilt:
\[
  d \in C^\interp{I}
  \;\iff\;
  \text{es existiert eine Subgraph-Einbettung }
  h: G_C \hookrightarrow G_\interp{I}
  \text{ mit } h(r_C) = d,
\]
wobei $G_C$ der \emph{Beschreibungsbaum} (description tree)
von $C$ ist und $r_C$ seine Wurzel.
\end{satz}

\begin{proof}
Wir beweisen beide Richtungen durch strukturelle Induktion
{\"u}ber den Aufbau von $C$.

\textbf{Definition der Einbettung.}
Der Beschreibungsbaum $G_C = (V_C, E_C, \lambda_C)$ von $C$
ist induktiv definiert als (vgl.\ Beweis von Satz~\ref{satz:lcs}):
\begin{itemize}
  \item $C = A$: einknötiger Baum, $V_C = \{r_C\}$,
        $\lambda_C(r_C) = \{A\}$, $E_C = \emptyset$.
  \item $C = \top$: einknötiger Baum, $\lambda_C(r_C) = \emptyset$.
  \item $C = C_1 \sqcap C_2$: Wurzel $r_C$ mit
        $\lambda_C(r_C) = \lambda_{C_1}(r_{C_1}) \cup \lambda_{C_2}(r_{C_2})$
        und den Teilb{\"a}umen von $C_1, C_2$ als Kinder (vereinigt).
  \item $C = \exists r.D$: Wurzel $r_C$ mit $r$-markierter
        Kante zu einem neuen Knoten $v_D$, der Wurzel des
        Beschreibungsbaums von $D$.
\end{itemize}

Eine \emph{Subgraph-Einbettung} $h: G_C \hookrightarrow G_\interp{I}$
ist ein injektiver Graphhomomorphismus, der:
\begin{enumerate}[label=(\alph*)]
  \item Knotenbeschriftungen respektiert:
        $\lambda_C(v) \subseteq \lambda(h(v))$ f{\"u}r alle $v \in V_C$,
  \item Kanten erhält:
        $(v, r, v') \in E_C \Rightarrow (h(v), h(v')) \in r^\interp{I}$.
\end{enumerate}

\textbf{($\Rightarrow$) Induktionsbeweis.}

\emph{Basisfall $C = A$:}
$d \in A^\interp{I}$ bedeutet $A \in \lambda(d)$.
Definiere $h(r_C) = d$.
Bedingung (a): $\lambda_C(r_C) = \{A\} \subseteq \lambda(d)$. \checkmark
Bedingung (b): $E_C = \emptyset$, nichts zu zeigen. \checkmark

\emph{Induktionsschritt $C = C_1 \sqcap C_2$:}
$d \in (C_1 \sqcap C_2)^\interp{I}$ bedeutet $d \in C_1^\interp{I}$
und $d \in C_2^\interp{I}$.
Nach Induktionshypothese existieren Einbettungen
$h_1: G_{C_1} \hookrightarrow G_\interp{I}$ mit $h_1(r_{C_1}) = d$
und $h_2: G_{C_2} \hookrightarrow G_\interp{I}$ mit $h_2(r_{C_2}) = d$.
Definiere $h = h_1 \cup h_2$ (beide Einbettungen werden vereinigt,
die Wurzel auf $d$ abgebildet). Bedingungen (a) und (b) gelten nach IH. \checkmark

\emph{Induktionsschritt $C = \exists r.D$:}
$d \in (\exists r.D)^\interp{I}$ bedeutet:
es gibt $e \in \Dom{I}$ mit $(d,e) \in r^\interp{I}$ und $e \in D^\interp{I}$.
Nach Induktionshypothese gibt es $h_D: G_D \hookrightarrow G_\interp{I}$
mit $h_D(r_D) = e$.
Definiere $h(r_C) = d$ und $h|_{V_D} = h_D$.
Bedingung (b) f{\"u}r die $r$-Kante $(r_C, r, r_D)$:
$(h(r_C), h(r_D)) = (d,e) \in r^\interp{I}$. \checkmark

\textbf{($\Leftarrow$) R{\"u}ckrichtung.}

Sei $h: G_C \hookrightarrow G_\interp{I}$ eine Einbettung mit $h(r_C) = d$.
Wir zeigen $d \in C^\interp{I}$ durch Induktion:

\emph{$C = A$:}
$\lambda_C(r_C) = \{A\} \subseteq \lambda(h(r_C)) = \lambda(d)$,
also $A \in \lambda(d)$, d.h.\ $d \in A^\interp{I}$. \checkmark

\emph{$C = C_1 \sqcap C_2$:}
Die Einbettung $h$ liefert $h|_{V_{C_1}}: G_{C_1} \hookrightarrow G_\interp{I}$
und $h|_{V_{C_2}}: G_{C_2} \hookrightarrow G_\interp{I}$, jeweils mit
Wurzelbild $d$. Per IH $d \in C_1^\interp{I}$ und $d \in C_2^\interp{I}$,
also $d \in (C_1 \sqcap C_2)^\interp{I}$. \checkmark

\emph{$C = \exists r.D$:}
Es gibt einen $r$-Nachfolger $v_D$ in $G_C$ von $r_C$.
Die Einbettung liefert $(h(r_C), h(v_D)) \in r^\interp{I}$.
Setze $e = h(v_D)$; dann $e \in D^\interp{I}$ nach IH.
Also $d \in (\exists r.D)^\interp{I}$. \checkmark
\qedhere
\end{proof}

\begin{figure}[ht]
  \centering
  \includegraphics[width=0.90\textwidth]{plots/plot7_dl_subgraph.pdf}
  \caption{Links: Eine DL-Interpretation als Rollengraph $G_\interp{I}$
           mit Knoten (Individuen) und beschrifteten Kanten (Rollen).
           Rechts: Das Konzept $\mathit{Doctor} \sqcap \exists\mathit{treats}.\mathit{Patient}$
           als Beschreibungsbaum $G_C$ -- seine Instanzierung durch \texttt{alice}
           entspricht genau einer Subgraph-Einbettung $G_C \hookrightarrow G_\interp{I}$.}
  \label{fig:dl-subgraph}
\end{figure}

\begin{korollar}[Subsumtion als Subgraph-Inklusion der Beschreibungsb{\"a}ume]
\label{kor:subsumtion-subgraph}
F{\"u}r \textsf{EL}-Konzepte $C, D$ ohne TBox gilt:
\[
  C \sqsubseteq D
  \;\iff\;
  \text{es existiert ein Homomorphismus } h: G_D \to G_C.
\]
Subsumtion ist {\"a}quivalent zur Existenz eines
\emph{Beschreibungsbaum-Homomorphismus} in umgekehrter Richtung.
\end{korollar}

\begin{proof}
Dies ist das klassische Resultat von Baader et al.~\cite{Baader2007}
(Theorem-Charakterisierung von \textsf{EL}-Subsumtion ohne TBox).
Ein Homomorphismus $h: G_D \to G_C$ existiert gdw.\ jede
Struktur, die $G_C$ als Subgraph enth{\"a}lt, auch $G_D$ enth{\"a}lt,
was {\"a}quivalent zu $C \sqsubseteq D$ ist.
Intuitiv: $D$ ist allgemeiner als $C$ gdw.\ der Beschreibungsbaum
von $D$ in den von $C$ eingebettet werden kann.
\end{proof}

\subsection{Signatur-Kodierung f{\"u}r ABox-Individuen}
\label{sec:signatur-abox}

\begin{definition}[ABox-Adjazenzmatrix und Signatur-Vektor]
\label{def:abox-sig}
Sei $\Abox$ eine ABox {\"u}ber Individuen $N_I = \{a_0, \ldots, a_{n-1}\}$
und Rollen $N_R = \{r_1, \ldots, r_m\}$.
F{\"u}r jede Rolle $r_k$ ist die \emph{$r_k$-Adjazenzmatrix}
$A^{(k)} \in \{0,1\}^{n \times n}$ definiert durch
$A^{(k)}_{ij} = 1$ gdw.\ $r_k(a_i, a_j) \in \Abox$.

Der \emph{Signatur-Vektor} von Individuum $a_j$
im Sinne des Subgraph-Algorithmus~\cite{Epp2026} ist:
\[
  \boldsymbol{\rho}(a_j) = \bigl(\rho_j^{(1)}, \ldots, \rho_j^{(m)}\bigr)
  \in \mathbb{N}^m,
  \quad
  \rho_j^{(k)} = \sum_{i=0}^{n-1} A^{(k)}_{ij} \cdot 2^i.
\]
\end{definition}

\begin{satz}[Signatur-Vektor kodiert die Rollennachbarschaft]
\label{satz:sig-nachbarschaft}
F{\"u}r Individuen $a_j, a_{j'} \in N_I$ gilt:
\begin{align*}
  &\boldsymbol{\rho}(a_j) = \boldsymbol{\rho}(a_{j'}) \;\iff \\
  &\quad\text{f{\"u}r alle } r_k \in N_R,\, a_i \in N_I:\;
   r_k(a_i, a_j) \in \Abox \;\Leftrightarrow\; r_k(a_i, a_{j'}) \in \Abox.
\end{align*}
Der Signatur-Vektor ist also ein eindeutiger Fingerabdruck der
eingehenden Rollennachbarschaft.
\end{satz}

\begin{proof}
\textbf{($\Rightarrow$):}
$\boldsymbol{\rho}(a_j) = \boldsymbol{\rho}(a_{j'})$ bedeutet
$\rho_j^{(k)} = \rho_{j'}^{(k)}$ f{\"u}r alle $k$.
Da $\rho_j^{(k)} = \sum_i A^{(k)}_{ij} \cdot 2^i$ die
Bin{\"a}rdarstellung des Spaltenvektors $(A^{(k)}_{0j}, \ldots, A^{(k)}_{(n-1)j})$
ist (Injektivit{\"a}t der Bin{\"a}rdarstellung), folgt
$A^{(k)}_{ij} = A^{(k)}_{ij'}$ f{\"u}r alle $i, k$.
Also $r_k(a_i, a_j) \in \Abox \Leftrightarrow r_k(a_i, a_{j'}) \in \Abox$.

\textbf{($\Leftarrow$):}
Wenn $A^{(k)}_{ij} = A^{(k)}_{ij'}$ f{\"u}r alle $i$,
dann sind die Spaltenvektoren identisch,
also $\rho_j^{(k)} = \rho_{j'}^{(k)}$ und somit
$\boldsymbol{\rho}(a_j) = \boldsymbol{\rho}(a_{j'})$.
\qedhere
\end{proof}

\begin{figure}[ht]
  \centering
  \includegraphics[width=0.90\textwidth]{plots/plot8_signatur_abox.pdf}
  \caption{Links: ABox-Adjazenzmatrix f{\"u}r eine medizinische Wissensbasis;
           die Spalten-Signaturen $\rho_j$ kodieren die eingehende
           Rollennachbarschaft jedes Individuums.
           Rechts: Der LCS zweier Individuen ergibt sich als
           ``Subgraph-Schnitt'' ihrer Signatur-Vektoren.}
  \label{fig:signatur-abox}
\end{figure}

\begin{satz}[Subgraph-LCS approximiert den semantischen LCS]
\label{satz:sublcs}
Seien $a_j, a_{j'}$ zwei ABox-Individuen mit MSCs
$\mathrm{MSC}(a_j)$ und $\mathrm{MSC}(a_{j'})$.
Der \emph{Subgraph-LCS} der Signatur-Vektoren,
definiert als:
\[
  \mathrm{SigLCS}(a_j, a_{j'})
  = \bigl(\min(\rho_j^{(k)}, \rho_{j'}^{(k)})\bigr)_{k=1}^{m},
\]
kodiert genau die gemeinsame Rollennachbarschaft beider Individuen
und ist eine untere Schranke f{\"u}r den semantischen LCS:
\[
  \mathrm{MSC}(a_j) \sqcap \mathrm{MSC}(a_{j'})
  \sqsubseteq
  \mathrm{lcs}(\mathrm{MSC}(a_j),\, \mathrm{MSC}(a_{j'})).
\]
\end{satz}

\begin{proof}
\textbf{Schritt 1 -- Semantik von SigLCS.}
Der Eintrag $\min(\rho_j^{(k)}, \rho_{j'}^{(k)})$ ist die
bitweise-AND-Verkn{\"u}pfung der Spalten-Vektoren:
$\min(A^{(k)}_{ij}, A^{(k)}_{ij'}) = A^{(k)}_{ij} \cdot A^{(k)}_{ij'}$.
Damit kodiert SigLCS$(a_j, a_{j'})$ genau jene Individuen $a_i$,
die \emph{beide} Individuen $a_j$ und $a_{j'}$ {\"u}ber Rolle $r_k$
als Vorg{\"a}nger haben:
\[
  r_k(a_i, a_j) \in \Abox \;\text{ und }\; r_k(a_i, a_{j'}) \in \Abox.
\]

\textbf{Schritt 2 -- Verbindung zum MSC.}
Nach Satz~\ref{satz:msc-exist} ist $\mathrm{MSC}(a_j)$
das spezifischste Konzept, das alle ABox-Assertionen
von $a_j$ widerspiegelt. Insbesondere enth{\"a}lt es den Konjunkt
$\exists r_k.\mathrm{MSC}(a_i)$ f{\"u}r alle $r_k(a_i, a_j) \in \Abox$.

\textbf{Schritt 3 -- Untere Schranke.}
Der LCS $\mathrm{lcs}(\mathrm{MSC}(a_j), \mathrm{MSC}(a_{j'}))$
ist das spezifischste Konzept, das von beiden MSCs subsumiert wird.
Da $\mathrm{MSC}(a_j) \sqcap \mathrm{MSC}(a_{j'})$ von beiden
subsumiert wird (Satz~\ref{satz:praeordnung}(iv)),
gilt per Definition des LCS:
\[
  \mathrm{MSC}(a_j) \sqcap \mathrm{MSC}(a_{j'})
  \sqsubseteq
  \mathrm{lcs}(\mathrm{MSC}(a_j), \mathrm{MSC}(a_{j'})).
\]

Der Subgraph-LCS entspricht dem Durchschnitt der Kantenmuster
und erzeugt damit einen Ausdruck, der von beiden MSCs
subsumiert wird -- er ist also eine untere Schranke.
\qedhere
\end{proof}

\subsection{Tableau-Blocking als Rotations{\"a}quivalenz}
\label{sec:blocking-rotation}

\begin{definition}[Signatur eines Tableau-Knotens]
\label{def:tableau-sig}
Sei $T = (V_T, E_T, \mathcal{L})$ ein \textsf{ALC}-Tableau.
Der \emph{Signatur-Vektor} eines Knotens $v \in V_T$ ist:
\[
  \boldsymbol{\sigma}(v)
  = \bigl(\,\rho_r \mid r \in N_R\bigr),
  \quad
  \rho_r = \sum_{A \in \mathcal{L}(v)} c_A \cdot 2^{i_A},
\]
wobei $c_A = 1$ falls $\exists r.A \in \mathcal{L}(v)$,
$i_A$ der Index von $A$ in einer festen Aufz{\"a}hlung
$\mathrm{sub}^+(C_0, \Tbox)$ ist.
\end{definition}

\begin{satz}[Blocking entspricht Signatur-{\"A}quivalenz]
\label{satz:blocking-aequiv}
Im Tableau-Algorithmus wird Knoten $v$ durch Vorfahre $u$ geblockt
(Pair-wise Blocking: $\mathcal{L}(u) = \mathcal{L}(v)$)
genau dann, wenn ihre Signatur-Vektoren gleich sind:
\[
  \mathcal{L}(u) = \mathcal{L}(v)
  \;\iff\;
  \boldsymbol{\sigma}(u) = \boldsymbol{\sigma}(v).
\]
\end{satz}

\begin{proof}
\textbf{($\Rightarrow$):}
$\mathcal{L}(u) = \mathcal{L}(v)$ bedeutet identische Mengen
von Konzeptausdr{\"u}cken.
Da der Signatur-Vektor $\boldsymbol{\sigma}(v)$ die Existenzformeln
$\exists r.A \in \mathcal{L}(v)$ als Bin{\"a}rzahl kodiert
und die Kodierung injektiv ist (nach Lemma 1 von~\cite{Epp2026}),
gilt $\boldsymbol{\sigma}(u) = \boldsymbol{\sigma}(v)$.

\textbf{($\Leftarrow$):}
$\boldsymbol{\sigma}(u) = \boldsymbol{\sigma}(v)$ bedeutet,
dass f{\"u}r alle Rollen $r$ und Konzepte $A$:
$\exists r.A \in \mathcal{L}(u) \Leftrightarrow \exists r.A \in \mathcal{L}(v)$.
Da die Tableau-Regeln nach vollst{\"a}ndiger Expansion alle
Boole'schen Konsequenzen aufgedeckt haben und das Tableau
vollst{\"a}ndig ist, bestimmt die Existenzstruktur die gesamte
Knotenbeschriftung. Also $\mathcal{L}(u) = \mathcal{L}(v)$.
\qedhere
\end{proof}

\begin{satz}[Tableau-Tiefe als Subgraph-Rotationsanzahl]
\label{satz:tableau-tiefe}
Die maximale Tiefe $d_{\max}$ des Tableau-Baums vor Blocking
stimmt mit der Anzahl der verschiedenen Signatur-Vektoren
{\"u}berein, die unter zyklischer Rotation von $\boldsymbol{\sigma}(v_0)$
entstehen:
\[
  d_{\max} \leq |\{\boldsymbol{\sigma}_k \mid k \in \{0,\ldots,n-1\}\}|,
\]
wobei $\boldsymbol{\sigma}_k$ die $k$-te zyklische Rotation ist.
\end{satz}

\begin{proof}
Nach Satz~\ref{satz:terminierung} wird Blocking ausgel{\"o}st,
sobald ein Knoten dieselbe Beschriftung wie ein Vorfahre hat.
Dies entspricht -- nach Satz~\ref{satz:blocking-aequiv} --
der {\"U}bereinstimmung der Signatur-Vektoren.
Da jede Tableau-Erweiterung eine neue Existenzformel aufdeckt und
damit den Signatur-Vektor ver{\"a}ndert, entspricht die Pfadl{\"a}nge
vor Blocking genau der Anzahl verschiedener erreichbarer
Signatur-Vektoren.
Im Subgraph-Algorithmus~\cite{Epp2026} werden genau $n$ zyklische
Rotationen gepr{\"u}ft (Lemma {\"u}ber Rotationsanzahl aus~\cite{Epp2026}),
was die obige Schranke liefert.
\end{proof}

\begin{figure}[ht]
  \centering
  \includegraphics[width=0.90\textwidth]{plots/plot9_tableau_graph.pdf}
  \caption{Tableau-Ast f{\"u}r ein \textsf{ALC}-Konzept als gerichteter Graph.
           Der Blocking-Knoten $v_4$ (orange) hat denselben Signatur-Vektor
           wie der Vorfahre $v_3$ -- dies entspricht exakt der
           Rotations{\"a}quivalenz im Subgraph-Algorithmus.}
  \label{fig:tableau-blocking}
\end{figure}

\begin{korollar}[Komplexit{\"a}tskonnex]
\label{kor:kompl-konnex}
Die Blocking-Tiefenschranke $2^{|C_0|}$ des Tableau-Algorithmus
(Satz~\ref{satz:terminierung}) entspricht der Anzahl verschiedener
Signatur-Vektoren -- also der maximalen Orbitgr{\"o}{\ss}e
$|\mathrm{Orb}(G)|$ im Sinne des Subgraph-Algorithmus.
Beide Gr{\"o}{\ss}en sind durch $2^{|\mathrm{sub}^+(C_0, \Tbox)|}$
beschr{\"a}nkt.
\end{korollar}

\begin{proof}
Nach Satz~\ref{satz:tableau-tiefe} gilt die Tiefenschranke durch
die Anzahl verschiedener Signatur-Vektoren.
Da jeder Signatur-Vektor eine Teilmenge von $\mathrm{sub}^+(C_0, \Tbox)$
repr{\"a}sentiert, gibt es h{\"o}chstens $2^{|\mathrm{sub}^+(C_0, \Tbox)|}$
verschiedene Vektoren. Die Blocking-Tiefe ist damit durch
$2^{|C_0| + |\Tbox|}$ beschr{\"a}nkt, was mit der
PSpace-Schranke aus Satz~\ref{satz:alc-pspace-upper}
{\"u}bereinstimmt.
\end{proof}

\subsection{Konjunktive Anfragen als Subgraph-Matching}
\label{sec:cq-matching}

\begin{definition}[Anfrage-Graph]
\label{def:anfragegraph}
Sei $q(\vec{x}) = \exists \vec{y}.\, \phi(\vec{x}, \vec{y})$ eine
konjunktive Anfrage (CQ, Definition~\ref{def:cert}).
Der \emph{Anfrage-Graph} $G_q = (V_q, E_q, \lambda_q)$ ist:
\begin{itemize}
  \item $V_q = \vec{x} \cup \vec{y}$ (alle Variablen),
  \item $(v, r, v') \in E_q$ gdw.\ $r(v, v')$ ein Konjunkt von $\phi$,
  \item $\lambda_q(v) = \{A \mid A(v) \text{ ein Konjunkt von } \phi\}$.
\end{itemize}
\end{definition}

\begin{satz}[CQ-Beantwortung als Subgraph-Einbettung]
\label{satz:cq-subgraph}
Sei $\KB = (\emptyset, \Abox)$ eine Wissensbasis ohne TBox,
$G_\Abox$ der Rollengraph der ABox.
F{\"u}r eine CQ $q(\vec{x})$ gilt:
\[
  \vec{a} \in \mathrm{cert}(q, \KB)
  \;\iff\;
  \text{es existiert eine Subgraph-Einbettung }
  h: G_q \hookrightarrow G_\Abox
  \text{ mit } h(\vec{x}) = \vec{a}.
\]
\end{satz}

\begin{proof}
\textbf{($\Rightarrow$):}
Sei $\vec{a} \in \mathrm{cert}(q, \KB)$.
Dann gilt $\Abox \models q(\vec{a})$ (da die einzige Interpretation
die kanonische ABox-Interpretation ist).
Also existiert eine Belegung $\theta: \vec{y} \to N_I$, so dass
alle Konjunkte $r(v,v')$ und $A(v)$ in der Belegung erf{\"u}llt sind.

Definiere $h: V_q \to N_I$ durch $h(x_i) = a_i$ und $h(y_j) = \theta(y_j)$.
Dann:
\begin{itemize}
  \item F{\"u}r $(v, r, v') \in E_q$: $r(h(v), h(v')) \in \Abox$,
        also $(h(v), h(v')) \in r^{\interp{I_\Abox}}$,
        also $(h(v), h(v')) \in E_\Abox$. \checkmark
  \item F{\"u}r $A \in \lambda_q(v)$: $A(h(v)) \in \Abox$,
        also $A \in \lambda(h(v))$. \checkmark
\end{itemize}
Also ist $h$ eine Subgraph-Einbettung $G_q \hookrightarrow G_\Abox$.

\textbf{($\Leftarrow$):}
Sei $h: G_q \hookrightarrow G_\Abox$ mit $h(\vec{x}) = \vec{a}$.
Setze $\theta(y_j) = h(y_j)$.
Dann ist $\theta$ eine Belegung, die alle Konjunkte erf{\"u}llt:
$r(h(v), h(v')) \in E_\Abox \Rightarrow r(h(v), h(v')) \in \Abox$
(da ABox-Graph und Rollengraph identisch sind).
Also $\Abox \models q(\vec{a})$ und damit $\vec{a} \in \mathrm{cert}(q, \KB)$.
\qedhere
\end{proof}

\begin{figure}[ht]
  \centering
  \includegraphics[width=0.90\textwidth]{plots/plot10_cq_subgraph.pdf}
  \caption{Links: Eine konjunktive Anfrage $q(x)$ als Anfrage-Graph $G_q$.
           Rechts: Die Einbettung $h: G_q \hookrightarrow G_\Abox$
           in den ABox-Rollengraph liefert die Antworte $\{alice\}$.
           CQ-Beantwortung ist strukturell {\"a}quivalent zum
           Subgraph-Matching-Problem.}
  \label{fig:cq-subgraph}
\end{figure}

\begin{satz}[Subgraph-Algorithmus f{\"u}r CQ-Beantwortung ohne TBox]
\label{satz:subgraph-cq}
F{\"u}r eine CQ $q$ und ABox $\Abox$ kann
$\mathrm{cert}(q, (\emptyset, \Abox))$ durch eine Variante des
Subgraph-Algorithmus~\cite{Epp2026} berechnet werden:
\begin{enumerate}[label=(\arabic*)]
  \item Berechne Signatur-Vektoren $\boldsymbol{\sigma}(G_q)$
        und $\boldsymbol{\sigma}(G_\Abox)$.
  \item F{\"u}r jede zyklische Rotation $k$ von $G_\Abox$:
        Pr{\"u}fe mittels LCS-Vergleich ob $G_q$ in der Rotation enthalten.
  \item Gib alle $a_i$ zur{\"u}ck, bei denen die Einbettung existiert.
\end{enumerate}
Die Laufzeit dieses Verfahrens betr{\"a}gt $O(|q|^2 \cdot |\Abox|^3)$.
\end{satz}

\begin{proof}
\textbf{Korrektheit:}
Nach Satz~\ref{satz:cq-subgraph} ist CQ-Beantwortung {\"a}quivalent zu
Subgraph-Einbettung. Der Subgraph-Algorithmus entscheidet genau
diese Frage -- korrekt nach dem Korrektheitssatz des
Subgraph-Algorithmus~\cite{Epp2026}.

\textbf{Laufzeit:}
Der Anfrage-Graph $G_q$ hat $|q|$ Knoten und Kanten.
Der ABox-Graph $G_\Abox$ hat $|\Abox|$ Knoten.
Signaturberechnung: $O(|q|^2)$ und $O(|\Abox|^2)$.
F{\"u}r jede der $|\Abox|$ Rotationen:
LCS-Vergleich in $O(|q| \cdot |\Abox|)$.
Gesamt: $O(|\Abox|^2) + O(|\Abox| \cdot |q| \cdot |\Abox|)
= O(|q| \cdot |\Abox|^2)$ pro Antwortkandidat,
und $|\Abox|$ Kandidaten ergeben $O(|q| \cdot |\Abox|^3)$.
\end{proof}

\begin{korollar}[Komplexit{\"a}tsvorteil f{\"u}r azyklische CQs]
\label{kor:azyklisch}
F{\"u}r \emph{azyklische} konjunktive Anfragen $q$ (azyklischer Anfrage-Graph)
gilt: Der Signatur-Vektor von $G_q$ hat Beschreibungstiefe $1$,
und der Subgraph-Algorithmus terminiert mit einer Rotation
in $O(|q| \cdot |\Abox|^2)$.
\end{korollar}

\begin{proof}
Bei azyklischem $G_q$ ist der Beschreibungsbaum flach
(Tiefe $\leq 1$), der Signatur-Vektor wird durch genau einen
Rotationsvergleich {\"u}berpr{\"u}ft, was den Faktor $|\Abox|$ spart.
\end{proof}

\begin{bemerkung}[Verbindung zu Datenbanktheorie]
Satz~\ref{satz:cq-subgraph} ist das klassische Resultat,
dass CQ-Auswertung {\"a}quivalent zu Homomorphismus-Existenz ist~\cite{Baader2007}.
Die neue Erkenntnis liegt in der \emph{algorithmischen Verbindung}:
Der Subgraph-Algorithmus~\cite{Epp2026} liefert f{\"u}r dieses Problem
einen konkreten Algorithmus mit expliziter Laufzeitanalyse,
der zudem unter SETH optimal ist.
\end{bemerkung}

% =======================================================================
\section{Zusammenfassung}
\label{sec:zusammenfassung}
% =======================================================================

Diese Arbeit hat Beschreibungslogiken von den mathematischen
Grundlagen bis zu neuartigen algorithmischen Verbindungen
vollst{\"a}ndig, formal und mit detaillierten Beweisen entwickelt.

\subsection{Ergebnisse im {\"U}berblick}

Die folgende Tabelle fasst alle in dieser Arbeit bewiesenen
Hauptresultate zusammen.

\begin{table}[ht]
\centering
\caption{{\"U}bersicht aller bewiesenen Hauptresultate}
\label{tab:zusammenfassung}
\begin{tabular}{@{}p{4.8cm}p{4.8cm}p{3cm}@{}}
\toprule
\textbf{Thema} & \textbf{Ergebnis} & \textbf{Nachweis} \\
\midrule
Signatur-Funktion & injektiv, injektiv (Lemma) & Abschn.~\ref{sec:grundlagen} \\
NNF-Transformation & poly.\ Zeit, Gr{\"o}{\ss}e $\leq 2|C|$ & Satz~\ref{satz:nnf} \\
Dualit{\"a}t $\exists$/$\forall$ & Mengengleichheit & Lemma~\ref{lem:duality} \\
\textsf{EL}-Normalisierung & poly.\ Zeit, $\leq 4|\Tbox|$ & Satz~\ref{satz:el-normal} \\
\textsf{EL}-Subsumtion & \textsc{P}-vollst{\"a}ndig, korrekt & Satz~\ref{satz:el-poly} \\
\textsf{ALC}-Erf{\"u}llbarkeit & \textsc{PSpace}-vollst{\"a}ndig & S{\"a}tze~\ref{satz:alc-pspace-upper},~\ref{satz:alc-pspace-lower} \\
\textsf{SHIQ} & \textsc{ExpTime}-vollst{\"a}ndig & Abschn.~\ref{sec:shiq-sroiq} \\
\textsf{SROIQ}/OWL~2 & \textsc{NExpTime}-vollst{\"a}ndig & Abschn.~\ref{sec:shiq-sroiq} \\
Tableau-Korrektheit & beide Richtungen, vollst{\"a}ndig & Satz~\ref{satz:tableau-korrektheit} \\
Tableau-Terminierung & Blocking-Tiefenschranke & Satz~\ref{satz:terminierung} \\
Subsumtion als Pr{\"a}ordnung & 5 Eigenschaften & Satz~\ref{satz:praeordnung} \\
Instanzpr{\"u}fung & Reduktion auf Unerf{\"u}llbarkeit & Satz~\ref{satz:instanz} \\
MSC in \textsf{EL} & existiert, poly.\ berechenbar & Satz~\ref{satz:msc-exist} \\
LCS in \textsf{EL} & Produktkonstruktion, poly.\ Gr{\"o}{\ss}e & Satz~\ref{satz:lcs} \\
Defeasible DL & Pr{\"a}ferenzmodell-Semantik & Abschn.~\ref{sec:nichtmonoton} \\
Repair-Existenz & garantiert & Satz~\ref{satz:repairs} \\
Brave Repair-Semantik & korrekt bei Konsistenz & Satz~\ref{satz:repair-korrektheit} \\
CQ-Datenkomplexit{\"a}t & \textsc{P} (\textsf{EL}) & Satz~\ref{satz:cq-data} \\
\midrule
\textbf{Neue Ergebnisse (Abschn.~\ref{sec:dl-subgraph})} & & \\
Instanz als Subgraph & Einbettungstheoretisch & Satz~\ref{satz:instanz-subgraph} \\
Subsumtion als Homom. & Beschreibungsbaum-Hom. & Korollar~\ref{kor:subsumtion-subgraph} \\
Signatur = Nachbarschaft & {\"A}quivalenz & Satz~\ref{satz:sig-nachbarschaft} \\
SigLCS untere Schranke & f{\"u}r semantischen LCS & Satz~\ref{satz:sublcs} \\
Blocking = Sig-{\"A}quiv. & f{\"u}r \textsf{ALC}-Tableau & Satz~\ref{satz:blocking-aequiv} \\
CQ = Subgraph-Matching & ohne TBox, exakt & Satz~\ref{satz:cq-subgraph} \\
\bottomrule
\end{tabular}
\end{table}

\subsection{Inhaltliche Zusammenfassung der Kapitel}

\textbf{Formale Grundlagen (Abschnitt~\ref{sec:grundlagen}).}
Die Arbeit beginnt mit einer vollst{\"a}ndigen Entwicklung der
Sprache und Semantik von \textsf{ALC} -- von der Signatur {\"u}ber
Interpretationen bis zu TBox, ABox und Wissensbasen.
Alle grundlegenden Eigenschaften (Dualit{\"a}t, Negationsnormalform,
Reduktion auf Unerf{\"u}llbarkeit) werden mit expliziten Beweisen belegt.

\textbf{Die Beschreibungslogik \textsf{EL} (Abschnitt~\ref{sec:el}).}
\textsf{EL} ist das Herzst{\"u}ck praktischer Ontologien (SNOMED~CT).
Wir beweisen, dass Subsumtion in \textsf{EL} in Polynomialzeit
entscheidbar ist (Satz~\ref{satz:el-poly}), und entwickeln
den vollst{\"a}ndigen Normalisierungs- und Regelanwendungsalgorithmus
mit Soundness- und Vollst{\"a}ndigkeitsbeweis {\"u}ber kanonisches Modell.

\textbf{Die Beschreibungslogik \textsf{ALC} (Abschnitt~\ref{sec:alc}).}
\textsf{ALC} ist die minimale boole'sch vollst{\"a}ndige DL.
Wir beweisen \textsc{PSpace}-Vollst{\"a}ndigkeit durch einen
nichtdeterministischen Tableau-Algorithmus (obere Schranke)
und explizite QBF-Reduktion (untere Schranke).

\textbf{Tableau-Verfahren (Abschnitt~\ref{sec:tableau}).}
Der Korrektheitsbeweis des Tableau-Verfahrens wird vollst{\"a}ndig
in beiden Richtungen -- Vollst{\"a}ndigkeit und Soundness --
gef{\"u}hrt, mit Induktion {\"u}ber alle sieben \textsf{ALC}-Konstruktoren.
Terminierung wird durch die Blocking-Tiefenschranke gesichert.

\textbf{Konzeptlernen (Abschnitt~\ref{sec:lernen}).}
MSC und LCS sind die zentralen Bausteine induktiven Lernens in DL.
Wir beweisen Existenz und Berechenbarkeit des MSC (Polynomialzeit-Konstruktion)
sowie des LCS durch vollst{\"a}ndige Produktkonstruktion auf Beschreibungsb{\"a}umen.

\textbf{Verbindung zum Subgraph-Algorithmus (Abschnitt~\ref{sec:dl-subgraph}).}
Das neue Kapitel dieser Arbeit stellt vier formale Verbindungen
zwischen Beschreibungslogik und dem Subgraph-Algorithmus~\cite{Epp2026} her:
DL-Interpretationen sind Rollengraphen; Konzept-Instanzierung ist
Subgraph-Einbettung; die Signatur-Funktion kodiert Rollennachbarschaften;
Tableau-Blocking entspricht Signatur-{\"A}quivalenz;
und CQ-Beantwortung ist {\"a}quivalent zu Subgraph-Matching.
Diese Verbindungen er{\"o}ffnen einen neuen algorithmischen Blickwinkel
auf DL-Reasoning.

% =======================================================================
\section{Ausblick}
\label{sec:ausblick}
% =======================================================================

Beschreibungslogiken stehen an der Schnittstelle von formaler Logik,
Wissensrepresentation, Algorithmentheorie und Datenbanktheorie.
Die in dieser Arbeit entwickelten Resultate -- insbesondere die
formalen Verbindungen zum Subgraph-Algorithmus -- er{\"o}ffnen
eine Vielzahl neuer Forschungsrichtungen.

\subsection{Subgraph-basiertes DL-Reasoning}

Die in Abschnitt~\ref{sec:dl-subgraph} hergestellten Verbindungen
zwischen DL und dem Subgraph-Algorithmus~\cite{Epp2026} deuten
auf ein bisher unerschlossenes algorithmisches Paradigma hin.

\begin{itemize}
  \item \textbf{Subgraph-Algorithmus als Reasoning-Backend.}
        Satz~\ref{satz:subgraph-cq} zeigt, dass CQ-Beantwortung
        ohne TBox direkt durch den Subgraph-Algorithmus l{\"o}sbar ist.
        \emph{Offene Frage:} L{\"a}sst sich der Subgraph-Algorithmus
        auf CQs mit TBox erweitern, etwa durch Rewriting der
        TBox-Konsequenzen in zus{\"a}tzliche ABox-Kanten?
        F{\"u}r \textsf{EL}-TBoxen ist dies durch den
        Erreichbarkeitsmengen-Algorithmus (Satz~\ref{satz:el-poly})
        m{\"o}glich: Man expandiert die ABox um alle
        durch die TBox erzwungenen Rollenfakten und wendet
        dann den Subgraph-Algorithmus an.

  \item \textbf{Rotationssymmetrie als Ontologie-Eigenschaft.}
        Korollar~\ref{kor:kompl-konnex} zeigt, dass die Tiefe
        des Tableau-Baums durch die Orbitgr{\"o}{\ss}e des Subgraph-Algorithmus
        beschr{\"a}nkt wird. Wenn eine Ontologie besonders viele
        rotationssymmetrische Konzepte enth{\"a}lt (gro{\ss}er Stabilisator),
        k{\"o}nnten wesentlich weniger Tableau-{\"A}ste exploriert werden.
        \emph{Offene Frage:} Welche Klassen von Ontologien besitzen
        systematische Rotationssymmetrien, und k{\"o}nnen diese
        f{\"u}r eine verbesserte Blocking-Strategie genutzt werden?

  \item \textbf{Optimale Laufzeit f{\"u}r strukturierte ABoxen.}
        F{\"u}r ABoxen mit zyklischer Struktur (Korollar~\ref{kor:azyklisch})
        kann CQ-Beantwortung effizienter durchgef{\"u}hrt werden.
        Allgemeiner: Die Orbits der Rotationsgruppe partitionieren
        den ABox-Graphen in {\"A}quivalenzklassen -- innerhalb einer
        Klasse muss CQ-Matching nur einmal durchgef{\"u}hrt werden.
        Dies k{\"o}nnte in gro{\ss}en, regelm{\"a}{\ss}ig strukturierten
        Wissensbasen erhebliche Laufzeiteinsparungen bringen.

  \item \textbf{Inkrementelles DL-Reasoning via Signatur-Updates.}
        Wenn eine ABox dynamisch um Rollenfakten erweitert wird,
        {\"a}ndert sich nach Satz~\ref{satz:sig-nachbarschaft}
        genau eine Signaturkomponente. Damit kann CQ-Matching
        inkrementell in $O(|q| \cdot |\Abox|^2)$ aktualisiert werden,
        statt in $O(|q| \cdot |\Abox|^3)$ neu berechnet werden zu m{\"u}ssen --
        ein Faktor $\Theta(|\Abox|)$ Speedup.
\end{itemize}

\subsection{Skalierbare Reasoning-Algorithmen}

Eine der dr{\"a}ngendsten praktischen Herausforderungen ist die Skalierung
von DL-Reasonern auf Ontologien mit Millionen von Konzepten und Fakten.
Aktuelle Reasoner wie HermiT~\cite{Glimm2014}, Pellet~\cite{Sirin2007}
und FaCT++~\cite{Tsarkov2006} sto{\ss}en bei sehr gro{\ss}en Ontologien
an empirische Grenzen.

\begin{itemize}
  \item \textbf{Modulare Ontologien.}
        Die Zerlegung gro{\ss}er Ontologien in kleinere Module~\cite{Grau2008}
        erlaubt unabh{\"a}ngige Verarbeitung.
        Formal ist ein $\Sigma$-Modul $M \subseteq \Tbox$ definiert durch:
        alle $\Sigma$-Konsequenzen von $\Tbox$ sind bereits Konsequenzen von $M$.
        F{\"u}r \textsf{EL} ist Modulextraktion in Polynomialzeit m{\"o}glich.
        Die Verbindung zum Subgraph-Algorithmus er{\"o}ffnet eine neue Perspektive:
        Module entsprechen zusammenh{\"a}ngenden Subgraphen des Rollengraphen.

  \item \textbf{Inkrementelles Reasoning.}
        Bei dynamischen Ontologien (Hinzuf{\"u}gen/Entfernen von Axiomen)
        soll die Erreichbarkeitsmenge $S(\cdot)$ (Satz~\ref{satz:el-poly})
        ohne vollst{\"a}ndige Neuberechnung aktualisiert werden~\cite{Baader2017}.
        Das Hinzuf{\"u}gen einer GCI $A \sqsubseteq B$ beeinflusst
        $S(X)$ genau f{\"u}r jene $X$ mit $A \in S(X)$ -- dies sind
        die Elemente des von $A$ aus erreichbaren Teilgraphen im
        Subgraph-Graphen der TBox.

  \item \textbf{Verteiltes und paralleles Reasoning.}
        Die $n$ Rotationsvergleiche des Subgraph-Algorithmus sind daten-parallel.
        F{\"u}r DL-Reasoning: Die CQ-Beantwortung f{\"u}r verschiedene
        Anfragekandidaten ist ebenfalls unabh{\"a}ngig parallelisierbar.
        Mit $p$ Prozessoren ergibt sich eine parallele Laufzeit von
        $O(|q| \cdot |\Abox|^3 / p + |\Abox|^2)$, was f{\"u}r gro{\ss}e
        ABoxen erhebliche Beschleunigung bedeutet.

  \item \textbf{Cache-optimierte Signatur-Berechnung.}
        Die Signatur-Vektoren der ABox-Individuen (Definition~\ref{def:abox-sig})
        k{\"o}nnen in einer kompakten Datenstruktur gespeichert werden,
        die sequentiell zugreifbar und daher cache-freundlich ist.
        Dies ist eine direkte algorithmische Konsequenz der
        Verbindung zum Subgraph-Algorithmus.
\end{itemize}

\subsection{Neuro-symbolische Integration}

\begin{itemize}
  \item \textbf{Subgraph-Einbettungen f{\"u}r DL-Konzepte.}
        Satz~\ref{satz:instanz-subgraph} {\"u}bersetzt Konzept-Instanzierung
        in Subgraph-Einbettungen. Dies liefert eine nat{\"u}rliche
        Grundlage f{\"u}r geometrische Einbettungen von DL-Konzepten:
        $d \in C^\interp{I}$ gdw.\ $\boldsymbol{\sigma}(G_C) \leq \boldsymbol{\sigma}(d)$
        (koordinatenweise). Diese Ordnungsrelation kann durch
        einfache neuronale Netzwerke approximiert werden.

  \item \textbf{LCS und MSC als Trainings-Ziele.}
        Der Subgraph-LCS (Satz~\ref{satz:sublcs}) liefert eine
        algorithmisch berechenbare Approximation des semantischen LCS.
        Neuronale Netze k{\"o}nnen trainiert werden, den LCS direkt
        aus Rollengraphen zu sch{\"a}tzen -- ohne expliziten MSC-Aufbau.
        Dies verbindet das Konzeptlernen (Abschnitt~\ref{sec:lernen})
        mit graph-neuronalen Netzwerken (GNNs).

  \item \textbf{Graphneuronale Netzwerke und Weisfeiler-Leman-Test.}
        Moderne GNNs sind in ihrer Ausdrucksst{\"a}rke durch den
        Weisfeiler-Leman-1-Test beschr{\"a}nkt.
        Die Signatur-Funktion des Subgraph-Algorithmus ist dem
        WL-1-Test {\"a}hnlich: Beide kodieren Nachbarschaftsstrukturen
        als Ganzzahlen. Damit entsteht die Frage, ob GNNs
        den Subgraph-Algorithmus und damit die CQ-Beantwortung
        (Satz~\ref{satz:cq-subgraph}) approximieren k{\"o}nnen.

  \item \textbf{Erkl{\"a}rbarkeit.}
        DL-Konzepte als Post-hoc-Erkl{\"a}rungen neuronaler Klassifikatoren
        k{\"o}nnen effizient {\"u}ber MSC und LCS berechnet werden
        (Abschnitt~\ref{sec:lernen}).
        Die Verbindung zum Subgraph-Algorithmus erm{\"o}glicht hier
        eine Skalierung auf sehr gro{\ss}e Wissensgraphen:
        Der Signatur-Vektor eines Individuums ist in $O(|\Abox|)$
        berechenbar und liefert sofort eine Aussage {\"u}ber
        die Konzeptzugeh{\"o}rigkeit.
\end{itemize}

\subsection{Erweiterungen der Ausdrucksst{\"a}rke}

\begin{itemize}
  \item \textbf{Subgraph-Algorithmus f{\"u}r gewichtete Rollenbeziehungen.}
        In \textsf{SHIQ} sind Zählrestriktionen $\geq n\,r.C$ erlaubt.
        Diese entsprechen gewichteten Kanten im Rollengraphen:
        Individuum $d$ hat $\geq n$ $r$-Nachfolger in $C^\interp{I}$.
        Die gewichtete Signatur-Erweiterung~\cite{Epp2026ext}
        kodiert diese Gewichte als $B$-adische Zahlen und
        erm{\"o}glicht effizienten Subgraph-Vergleich mit Zählbedingungen.

  \item \textbf{Temporale DL und dynamische Subgraphen.}
        In $\mathcal{ALC}_{LTL}$~\cite{Artale2001} ver{\"a}ndern sich
        Konzepterweiterungen {\"u}ber die Zeit.
        Dies entspricht einem dynamischen Rollengraphen, in dem
        Kanten hinzugef{\"u}gt und entfernt werden.
        Der inkrementelle Signatur-Update (Satz~\ref{satz:lokaler-update}
        aus~\cite{Epp2026ext}) liefert hier einen $O(n^2)$-Algorithmus
        f{\"u}r temporale CQ-Beantwortung.

  \item \textbf{Fuzzy-DL und gewichtete Subgraph-Einbettungen.}
        In Fuzzy-DL~\cite{Straccia2001} sind Konzeptzugeh{\"o}rigkeiten
        Werte in $[0,1]$.
        Der Subgraph-Algorithmus mit gewichteten Kanten
        (rationale Gewichte) liefert nat{\"u}rlich eine
        graduelle Subgraph-Einbettung: $d$ geh{\"o}rt zu Grad
        $\min_{(v,r,v') \in G_C} w(h(v), h(v'))$ dem Konzept $C$ an.

  \item \textbf{EL mit konkreten Dom{\"a}nen und Signatur-Erweiterung.}
        Werterestriktionen $\exists r.\{x \mid x \geq t\}$ in
        $\mathcal{EL}(\mathcal{D})$ lassen sich in die gewichtete
        Signatur einbetten: Kante $(a_i, a_j)$ mit Gewicht $w_{ij} = w(a_j)$,
        Zählbedingung $w_{ij} \geq t$.
        Damit wird CQ-Beantwortung mit arithmetischen Einschr{\"a}nkungen
        auf gewichtetes Subgraph-Matching reduziert.
\end{itemize}

\subsection{Lerntheorie und Konzeptlernen}

\begin{itemize}
  \item \textbf{PAC-Lernbarkeit und Subgraph-Komplexit{\"a}t.}
        PAC-Lernbarkeit~\cite{Valiant1984} erfordert polynomielle
        Stichprobenkomplexit{\"a}t. F{\"u}r \textsf{EL}-Konzepte der
        Rollentiefe $\leq k$ ist PAC-Lernbarkeit bekannt.
        Die Verbindung zum Subgraph-Algorithmus liefert eine
        neue Charakterisierung: Ein \textsf{EL}-Konzept der Tiefe $k$
        entspricht einem Beschreibungsbaum der Tiefe $k$,
        der als Subgraph-Template mit $k$ Rotationsschritten verglichen wird.
        Je kleiner $k$, desto weniger Rotationen werden ben{\"o}tigt,
        desto schneller ist das Lernen.

  \item \textbf{Untere Schranken durch SETH.}
        Unter SETH~\cite{Epp2026} ist der Subgraph-Algorithmus
        optimal in $\Theta(n^3)$. Dies {\"u}bertr{\"a}gt sich auf
        CQ-Beantwortung: F{\"u}r allgemeine CQs kann man unter SETH
        keine wesentlich schnellere Methode erwarten.
        F{\"u}r strukturierte ABoxen (azyklisch, beschr{\"a}nkte Treewidth)
        gibt es jedoch Ausnahmen (Korollar~\ref{kor:azyklisch}).

  \item \textbf{Lernen unter Inkonsistenz.}
        Wenn die ABox Fehler enth{\"a}lt (inkonsistente Rollenfakten),
        liefert der Subgraph-Algorithmus trotzdem eine maximale
        Subgraph-Einbettung -- analog zur Brave-Semantik der Repairs
        (Satz~\ref{satz:repair-korrektheit}).
        Dies verbindet repair-basiertes Reasoning mit
        robusten Matching-Verfahren.

  \item \textbf{Aktives Lernen mit Subgraph-Anfragen.}
        In aktivem Lernen stellt der Algorithmus gezielt Anfragen.
        Eine nat{\"u}rliche Anfragestrategie: Frage, ob ein bestimmter
        Subgraph $G_q$ in $G_\Abox$ eingebettet werden kann.
        Die Antwort -- ja oder nein -- verfeinert die Hypothese.
        Die optimale Anzahl solcher Anfragen ist durch die
        LCS-Gr{\"o}{\ss}e beschr{\"a}nkt (Satz~\ref{satz:lcs}).
\end{itemize}

\subsection{Linked Data, SPARQL und Wissensgraphen}

\begin{itemize}
  \item \textbf{Subgraph-Algorithmus f{\"u}r SPARQL.}
        SPARQL-Anfragen sind strukturell CQs (mit optionalen und
        Unionskonstrukten). F{\"u}r den konjunktiven Kern gilt
        Satz~\ref{satz:cq-subgraph}: SPARQL-Anfrageauswertung
        {\"u}ber einer ABox ist Subgraph-Matching.
        Der Subgraph-Algorithmus liefert damit eine
        konkurrenzf{\"a}hige Alternative zu Standard-SPARQL-Engines
        f{\"u}r strukturierte Anfragen.

  \item \textbf{Wissensgraph-Vervollst{\"a}ndigung.}
        Gro{\ss}e Wissensgraphen (Wikidata, DBpedia) sind unvollst{\"a}ndig.
        Ein fehlender Rollenfakt $r(a, b)$ {\"a}ndert den Signatur-Vektor
        von $b$ um genau $2^{\mathrm{idx}(a)}$.
        Durch Vergleich mit bekannten vollst{\"a}ndigen Subgraphen
        (Templates) k{\"o}nnen fehlende Fakten als ``Subgraph-L{\"u}cken''
        identifiziert und erg{\"a}nzt werden.

  \item \textbf{Ontology-Mediated Query Answering.}
        F{\"u}r OWL~2~QL ist Anfragebeantwortung in \textsc{AC}$^0$~\cite{Calvanese2007}.
        Die algorithmische Br{\"u}cke zum Subgraph-Algorithmus er{\"o}ffnet
        die Frage, ob Subgraph-Techniken f{\"u}r OWL~2~EL
        (\textsc{P}-vollst{\"a}ndig in Datenkomplexit{\"a}t, Satz~\ref{satz:cq-data})
        praktisch effizientere Implementierungen erm{\"o}glichen.
\end{itemize}

\subsection{Verbindungen zur Theoretischen Informatik}

\begin{itemize}
  \item \textbf{Homomorphismus-Theorie.}
        Korollar~\ref{kor:subsumtion-subgraph} und Satz~\ref{satz:cq-subgraph}
        verbinden DL-Reasoning mit der Homomorphismus-Theorie
        (vgl.\ CSP -- Constraint Satisfaction Problems).
        Die Frage, ob CQ-Beantwortung mit TBox auf allgemeine
        Homomorphismusprobleme reduzierbar ist, ber{\"u}hrt
        das Dichotomie-Theorem von Bulatov~\cite{Bulatov2017}:
        Entweder ist ein CSP-Problem in \textsc{P} oder \textsc{NP}-vollst{\"a}ndig.
        F{\"u}r strukturierte ABoxen (feste Treewidth) f{\"a}llt
        CQ-Beantwortung in \textsc{P}.

  \item \textbf{Bisimulation und Rotations{\"a}quivalenz.}
        Bisimulante Elemente erf{\"u}llen dieselben \textsf{ALC}-Konzepte~\cite{Blackburn2001}.
        Satz~\ref{satz:blocking-aequiv} zeigt, dass Blocking im
        Tableau {\"a}quivalent zu Signatur-{\"A}quivalenz ist.
        Die Frage, ob Bisimulation durch Rotations{\"a}quivalenz der
        Signatur-Vektoren charakterisiert werden kann, verbindet
        DL-Semantik mit dem Subgraph-Algorithmus auf fundamentale Weise.

  \item \textbf{Fixpunktlogiken und iterativer Subgraph-Vergleich.}
        Der $\mu$-Kalk{\"u}l erm{\"o}glicht Transitivit{\"a}t.
        Im Subgraph-Kontext entspricht der transitive Abschluss
        der iterativen Anwendung des Signatur-Vergleichs:
        Starte mit $G_q$, f{\"u}ge sukzessiv $r$-Nachfolger hinzu
        und pr{\"u}fe nach jedem Schritt auf Subgraph-{\"U}bereinstimmung.
        Dies ist {\"a}quivalent zum Least-Fixpoint-Reasoning im $\mu$-Kalk{\"u}l.

  \item \textbf{Guardierte Logiken.}
        \textsf{ALC} ist im guardierten Fragment (GF) der
        Pr{\"a}dikatenlogik eingebettet~\cite{Gradel1999}.
        Die Subgraph-Signatur beschr{\"a}nkt implizit die m{\"o}glichen
        Unifikationen: Nur ``bewachte'' Paare $(d, e)$ mit
        gemeinsamer guardierender Kante k{\"o}nnen unifiziert werden.
        Diese Guard-Bedingung entspricht einer Einschr{\"a}nkung
        der Subgraph-Einbettungen auf lokale Umgebungen.
\end{itemize}

\subsection{Anwendungen und Standardisierung}

\begin{itemize}
  \item \textbf{Biomedizin: SNOMED~CT und subgraph-basiertes Reasoning.}
        SNOMED~CT~\cite{SNOMED} modelliert {\"u}ber 350.000 klinische
        Konzepte in \textsf{EL} mit komplexer Rollenstruktur.
        Die Verbindung zum Subgraph-Algorithmus er{\"o}ffnet die M{\"o}glichkeit,
        {\"a}hnliche Konzepte durch Subgraph-Distanz (metrische Eigenschaften)
        zu finden -- dies w{\"a}re f{\"u}r medizinische Suche n{\"u}tzlich.

  \item \textbf{Rechtliche Wissensmodellierung.}
        Gesetze sind formal als DL-Ontologien kodierbar.
        Gesetzliche {\"A}hnlichkeit (``Gesetz A {\"a}hnelt Gesetz B'') ist
        formalisierbar als Subgraph-Distanz der Rollengraphen beider Gesetze.

  \item \textbf{Autonome Systeme und Szenenverst{\"a}ndnis.}
        Szenenbeschreibungen (Objekte, Relationen) entsprechen
        nat{\"u}rlich Rollengraphen.
        Konzept-Matching (``Ist das ein Fahrzeug vor mir?'')
        entspricht Subgraph-Einbettung des Fahrzeug-Konzeptgraphen
        in den Szenen-Rollengraphen.
        Dies verbindet DL-Ontologien mit computer-vision-basierten
        Szenengraphen.
\end{itemize}

\subsection{Offene Theoretische Fragen}

\begin{enumerate}
  \item \textbf{Subgraph-Algorithmus mit TBox-Konsequenzen.}
        Kann der Subgraph-Algorithmus so erweitert werden,
        dass er bei der Signaturberechnung automatisch
        TBox-Ableitungen ber{\"u}cksichtigt?
        F{\"u}r \textsf{EL} ist dies durch Expandierung der ABox via
        $S(\cdot)$-Berechnung konzeptionell m{\"o}glich,
        formal aber noch nicht ausgearbeitet.

  \item \textbf{Optimale Blocking-Strategie via Rotations-Symmetrien.}
        Gibt es eine Blocking-Strategie f{\"u}r \textsf{ALC},
        die Rotationssymmetrien der Tableau-Signatur ausnutzt
        und damit exponentiell weniger Bl{\"o}cke als Pair-wise Blocking ben{\"o}tigt?

  \item \textbf{LCS und Subgraph-LCS -- exakte Charakterisierung.}
        Satz~\ref{satz:sublcs} gibt eine untere Schranke.
        Wann gilt exakte Gleichheit: $\mathrm{SigLCS} = \mathrm{lcs}$?
        Eine vollst{\"a}ndige Charakterisierung der ABoxen, f{\"u}r die
        der Subgraph-LCS den semantischen LCS exakt trifft,
        ist ein offenes Problem.

  \item \textbf{Entscheidbarkeitsgrenze f{\"u}r Subgraph-Erweiterungen.}
        Komplexe Rollenaxiome in \textsf{SROIQ} f{\"u}hren zur
        \textsc{NExpTime}-Vollst{\"a}ndigkeit. Im Subgraph-Kontext
        entsprechen diese Rollenkompositions-Einschr{\"a}nkungen auf
        den Subgraphen der ABox. Welche Einschr{\"a}nkungen sind
        gerade noch entscheidbar?

  \item \textbf{Verbindung zur Graphenisomorphie.}
        Graphenisomorphie (GI) liegt in Quasipolynomialzeit~\cite{Babai2016}.
        Der Subgraph-Algorithmus l{\"o}st eine eingeschr{\"a}nktere Frage
        (zyklische Rotationen) in $O(n^3)$.
        Kann man durch mehrfache Anwendung des Subgraph-Algorithmus
        auf Teilgraphen einen GI-Algorithmus n{\"a}hern,
        und wie verh{\"a}lt sich dessen Komplexit{\"a}t zur Quasi\-polynomialzeit-Schranke?
\end{enumerate}

\subsection{Abschlie{\ss}ende Einordnung}

Beschreibungslogiken verbinden rigorose formale Semantik mit
effizienten Algorithmen und vielf{\"a}ltigen Anwendungen.
Diese Arbeit hat gezeigt, dass DL nicht isoliert steht,
sondern tiefe strukturelle Verbindungen zu klassischen
Algorithmen-Problemen -- insbesondere dem Subgraph-Problem --
aufweist.

Die in Abschnitt~\ref{sec:dl-subgraph} entwickelten Verbindungen
sind nicht blo{\ss} Analogien, sondern formale {\"A}quivalenzen:
\begin{itemize}
  \item Konzept-Instanzierung \emph{ist} Subgraph-Einbettung
        (Satz~\ref{satz:instanz-subgraph}).
  \item Signatur-Vektoren \emph{sind} Rollennachbarschafts-Kodierungen
        (Satz~\ref{satz:sig-nachbarschaft}).
  \item Tableau-Blocking \emph{ist} Signatur-{\"A}quivalenz
        (Satz~\ref{satz:blocking-aequiv}).
  \item CQ-Beantwortung \emph{ist} Subgraph-Matching
        (Satz~\ref{satz:cq-subgraph}).
\end{itemize}

Diese Gleichsetzungen er{\"o}ffnen eine neue algorithmische Perspektive:
Der Subgraph-Algorithmus~\cite{Epp2026} mit seiner optimalen
$O(n^3)$-Laufzeit (unter SETH) liefert ein konkretes Werkzeug,
das direkt f{\"u}r DL-Reasoning eingesetzt werden kann.
Umgekehrt liefert die reiche Theorie der Beschreibungslogiken --
Halbordnungen, LCS, Repair-Semantik, Bisimulation --
theoretische Fundierungen f{\"u}r den Subgraph-Algorithmus.

Die vier gro{\ss}en Forschungsrichtungen der Zukunft
an dieser Schnittstelle sind:
\begin{enumerate}
  \item \textbf{Subgraph-basiertes Reasoning:}
        Einsatz des Subgraph-Algorithmus als effizienter
        Reasoning-Backend f{\"u}r CQ-Beantwortung in gro{\ss}en ABoxen.
  \item \textbf{Skalierung:}
        Nutzung der inkrementellen Signatur-Updates f{\"u}r
        dynamische Ontologien und Echtzeit-Anfragen.
  \item \textbf{Integration:}
        Verbindung von Subgraph-Einbettungen mit graph-neuronalen
        Netzwerken f{\"u}r neuro-symbolisches Reasoning.
  \item \textbf{Theorie:}
        Formale Charakterisierung der Grenzen dieser
        {\"A}quivalenzen (mit TBox, ausdrucksstarken DL-Fragmenten,
        unter Standardhypothesen wie SETH).
\end{enumerate}

Der Kurs L.079.05832 ``Introduction to Description Logics''
an der Universit{\"a}t Bielefeld, geleitet von
Prof.\ Dr.\ Anni-Yasmin Turhan, vermittelt die in dieser Arbeit
formalisierte Grundlage. Die hier hergestellten Verbindungen
zum Subgraph-Algorithmus~\cite{Epp2026} er{\"o}ffnen ein neues,
vielversprechendes Forschungsfeld an der Schnittmenge von
Beschreibungslogik, Graphentheorie und Algorithmentheorie.

% =======================================================================
\begin{thebibliography}{99}
\bibitem{Babai2016}
L.~Babai.
Graph Isomorphism in Quasipolynomial Time.
In \textit{Proc.\ ACM STOC 2016}, Seiten 684--697, 2016.

\bibitem{Bulatov2017}
A.~Bulatov.
A Dichotomy Theorem for Nonuniform CSPs.
In \textit{Proc.\ IEEE FOCS 2017}, Seiten 319--330, 2017.

\bibitem{Epp2026}
S.~Epp.
\textit{Der Subgraph-Algorithmus}.
Technischer Bericht, Universit{\"{a}}t Bielefeld, 2026.
\url{https://gitlab.com/epp-group}

\bibitem{Epp2026ext}
S.~Epp.
\textit{Neue Erkenntnisse zum Subgraph-Algorithmus}.
Technischer Bericht, Universit{\"{a}}t Bielefeld, 2026.



\bibitem{Artale2001}
A.~Artale und E.~Franconi.
Temporal Description Logics.
In \textit{Handbook of Temporal Reasoning in Artificial Intelligence},
Elsevier, 2001.

\bibitem{Baader1999}
F.~Baader.
Computing the Least Common Subsumer in the Description Logic
$\mathcal{EL}$ w.r.t.\ Terminological Cycles.
\textit{LTCS-Report}, TU Dresden, 1999.

\bibitem{Baader2005}
F.~Baader, S.~Brandt und C.~Lutz.
Pushing the $\mathcal{EL}$ Envelope.
In \textit{Proc.\ IJCAI 2005}, Seiten 364--369, 2005.

\bibitem{Baader2007}
F.~Baader, D.~Calvanese, D.~McGuinness, D.~Nardi und P.~Patel-Schneider (Hrsg.).
\textit{The Description Logic Handbook: Theory, Implementation
and Applications}, 2.\ Aufl.
Cambridge University Press, 2007.

\bibitem{Baader2017}
F.~Baader, I.~Horrocks, C.~Lutz und U.~Sattler.
\textit{An Introduction to Description Logic}.
Cambridge University Press, 2017.

\bibitem{Bienvenu2014}
M.~Bienvenu und M.~Ortiz.
Ontology-Mediated Query Answering with Data-Tractable Description Logics.
In \textit{Proc.\ Reasoning Web Summer School},
LNCS 8714, 2014.

\bibitem{Blackburn2001}
P.~Blackburn, M.~de~Rijke und Y.~Venema.
\textit{Modal Logic}.
Cambridge University Press, 2001.

\bibitem{Bonatti2015}
P.~A.~Bonatti, M.~Faella und L.~Sauro.
Defeasible Inclusions in Low-Complexity DLs.
\textit{Journal of Artificial Intelligence Research},
42:719--764, 2015.

\bibitem{Brachman1979}
R.~J.~Brachman.
On the Epistemological Status of Semantic Networks.
In \textit{Associative Networks}, Academic Press, 1979.

\bibitem{Brachman1985}
R.~J.~Brachman und H.~J.~Levesque.
The Tractability of Subsumption in Frame-Based
Description Languages.
In \textit{Proc.\ AAAI-84}, Seiten 34--37, 1985.

\bibitem{Britz2021}
K.~Britz und G.~Casini.
A KLM Perspective on Defeasible Reasoning for Description Logics.
\textit{IfCoLog Journal of Logics and their Applications},
8(7), 2021.

\bibitem{Buchanan1984}
B.~G.~Buchanan und E.~H.~Shortliffe.
\textit{Rule-Based Expert Systems: The MYCIN Experiments}.
Addison-Wesley, 1984.

\bibitem{Calvanese2007}
D.~Calvanese, G.~De~Giacomo, D.~Lembo, M.~Lenzerini und R.~Rosati.
Tractable Reasoning and Efficient Query Answering in Description Logics:
The DL-Lite Family.
\textit{Journal of Automated Reasoning}, 39(3):385--429, 2007.

\bibitem{Chen2021}
J.~Chen, H.~Hu, E.~Jimenez-Ruiz, O.~Munn und I.~Horrocks.
OWL2Vec*: Embedding of OWL Ontologies.
\textit{Machine Learning}, 110(7):1813--1845, 2021.

\bibitem{Donini1997}
F.~M.~Donini, M.~Lenzerini, D.~Nardi und A.~Schaerf.
Reasoning in Description Logics.
In \textit{Principles of Knowledge Representation and Reasoning},
CSLI Publications, 1997.

\bibitem{Glimm2014}
B.~Glimm, I.~Horrocks, B.~Motik, G.~Stoilos und Z.~Wang.
HermiT: An OWL 2 Reasoner.
\textit{Journal of Automated Reasoning}, 53(3):245--269, 2014.

\bibitem{Gradel1999}
E.~Gradel.
On the Restraining Power of Guards.
\textit{Journal of Symbolic Logic}, 64(4):1719--1742, 1999.

\bibitem{Grau2008}
B.~C.~Grau, I.~Horrocks, Y.~Kazakov und U.~Sattler.
Modular Reuse of Ontologies: Theory and Practice.
\textit{Journal of Artificial Intelligence Research},
31:273--318, 2008.

\bibitem{Hohenecker2020}
P.~Hohenecker und T.~Lukasiewicz.
Ontology Reasoning with Deep Neural Networks.
\textit{Journal of Artificial Intelligence Research},
68:503--540, 2020.

\bibitem{Horrocks2000}
I.~Horrocks und U.~Sattler.
Ontology Reasoning in the SHIQ(D) Description Logic.
In \textit{Proc.\ IJCAI 2001}, Seiten 199--204, 2001.

\bibitem{Horrocks2006}
I.~Horrocks, O.~Kutz und U.~Sattler.
The Even More Irresistible SROIQ.
In \textit{Proc.\ KR 2006}, Seiten 57--67, 2006.

\bibitem{Kulmanov2019}
M.~Kulmanov, W.~Liu-Wei, Y.~Yan und R.~Hoehndorf.
EL Embeddings: Geometric Construction of Models for
the Description Logic EL++.
In \textit{Proc.\ IJCAI 2019}, Seiten 6103--6109, 2019.

\bibitem{Lehmann2009}
J.~Lehmann und P.~Hitzler.
Concept Learning in Description Logics Using Refinement Operators.
\textit{Machine Learning}, 78(1--2):203--250, 2010.

\bibitem{McCarthy1963}
J.~McCarthy und P.~J.~Hayes.
Some Philosophical Problems from the Standpoint of
Artificial Intelligence.
In \textit{Machine Intelligence 4}, Edinburgh University Press, 1969.

\bibitem{Minsky1975}
M.~Minsky.
A Framework for Representing Knowledge.
\textit{MIT AI Memo 306}, 1974.

\bibitem{OWL1}
D.~L.~McGuinness und F.~van~Harmelen.
OWL Web Ontology Language Overview.
\textit{W3C Recommendation}, 2004.
\url{https://www.w3.org/TR/owl-features/}

\bibitem{OWL2Primer}
P.~Hitzler, M.~Krotzsch, B.~Parsia, P.~F.~Patel-Schneider
und S.~Rudolph.
OWL 2 Web Ontology Language Primer.
\textit{W3C Recommendation}, 2009.
\url{https://www.w3.org/TR/owl2-primer/}

\bibitem{Pan2006}
J.~Z.~Pan und E.~Thomas.
Approximating OWL-DL Ontologies.
In \textit{Proc.\ AAAI 2007}, Seiten 1434--1439, 2007.

\bibitem{Penaloza2017}
R.~Penaloza und B.~Sertkaya.
Understanding the Complexity of Axiom Pinpointing
in Light-Weight DLs.
\textit{Artificial Intelligence}, 250:80--104, 2017.

\bibitem{Quillian1968}
M.~R.~Quillian.
Semantic Memory.
In \textit{Semantic Information Processing},
MIT Press, 1968.

\bibitem{Savitch1970}
W.~J.~Savitch.
Relationships Between Nondeterministic and Deterministic
Tape Complexities.
\textit{Journal of Computer and System Sciences},
4(2):177--192, 1970.

\bibitem{Schmidt1991}
K.~Schild.
A Correspondence Theory for Terminological Logics:
Preliminary Report.
In \textit{Proc.\ IJCAI 1991}, Seiten 466--471, 1991.

\bibitem{Sirin2007}
E.~Sirin, B.~Parsia, B.~C.~Grau, A.~Kalyanpur und Y.~Katz.
Pellet: A Practical OWL-DL Reasoner.
\textit{Journal of Web Semantics}, 5(2):51--53, 2007.

\bibitem{SNOMED}
SNOMED International.
\textit{SNOMED CT -- Systematized Nomenclature of Medicine
Clinical Terms}.
\url{https://www.snomed.org}, 2024.

\bibitem{Stockmeyer1976}
L.~J.~Stockmeyer.
The Polynomial-Time Hierarchy.
\textit{Theoretical Computer Science}, 3(1):1--22, 1976.

\bibitem{Straccia2001}
U.~Straccia.
Reasoning within Fuzzy Description Logics.
\textit{Journal of Artificial Intelligence Research},
14:137--166, 2001.

\bibitem{Tobies2001}
S.~Tobies.
\textit{Complexity Results and Practical Algorithms for
Logics in Knowledge Representation}.
Dissertation, RWTH Aachen, 2001.

\bibitem{Tsarkov2006}
D.~Tsarkov und I.~Horrocks.
FaCT++ Description Logic Reasoner: System Description.
In \textit{Proc.\ IJCAR 2006}, LNAI 4130, Seiten 292--297, 2006.

\bibitem{Turhan2012}
A.-Y.~Turhan.
Learnability of Description Logic Programs.
In \textit{Proc.\ ISWC 2012}, 2012.

\bibitem{Valiant1984}
L.~G.~Valiant.
A Theory of the Learnable.
\textit{Communications of the ACM}, 27(11):1134--1142, 1984.

\end{thebibliography}

\end{document}
